% Section file for Chapter: Twisted Holography and Quantum Gravity
% Result (G10): Fredholm Partition Functions
%
% This section develops the nonperturbative partition function
% as a Fredholm determinant on the one-particle Bergman space,
% extending from the Heisenberg prototype through the full
% standard landscape.

% ==========================================
% LOCAL MACRO PROVISIONS
% ==========================================
\providecommand{\Bergman}{\mathcal{B}}
\providecommand{\Fock}{\mathcal{F}}
\providecommand{\Hilb}{\mathcal{H}\!\mathit{ilb}}
\providecommand{\sewop}{T}
\providecommand{\sewker}{K}
\providecommand{\trcl}{\mathcal{L}^1}
\providecommand{\hscl}{\mathcal{L}^2}
\providecommand{\freddet}{\mathrm{det}_{\mathrm{Fred}}}
\providecommand{\Wick}{\mathcal{W}}
\providecommand{\Bergker}{\mathsf{B}}
\providecommand{\primeform}{E}
\providecommand{\collar}{\mathfrak{c}}
\providecommand{\sewenv}{\cA^{\mathrm{sew}}}
\providecommand{\anbar}{\barB^{\mathrm{an}}}
\providecommand{\algcore}{\cA_{\mathrm{alg}}}
\providecommand{\codercat}{D^{\mathrm{co}}}
\providecommand{\contrcat}{D^{\mathrm{ctr}}}
\providecommand{\graphnorm}{\|\cdot\|_{\mathrm{graph}}}
\providecommand{\secquant}{\Gamma}
\providecommand{\hodgebdle}{\mathbb{E}}
\providecommand{\hodgedet}{\det\hodgebdle}
\providecommand{\zetareg}{\det\nolimits'_\zeta}
\providecommand{\dedeta}{\eta}
\providecommand{\Sp}{\mathrm{Sp}}
\providecommand{\Siegel}{\mathfrak{H}}
\providecommand{\OPEcoeff}{C}

\section{Fredholm partition functions}
\label{sec:thqg-fredholm-partition-functions}
\index{Fredholm determinant!partition function|textbf}
\index{partition function!Fredholm}

The preceding nine results of this chapter extract holographic,
gravitational, and bootstrap consequences from the universal
Maurer--Cartan element
$\Theta_\cA \in \MC(\gAmod)$
at the level of formal algebraic identities.  This final result
crosses from the algebraic to the analytic: the partition functions
that witness the genus expansion are not merely formal power series
but convergent Fredholm determinants on explicit Hilbert spaces.

The key bridge is the \emph{sewing envelope}
$\sewenv$
(Definition~\ref{def:sewing-envelope}): the Hausdorff completion
of the algebraic chiral core $\algcore$ for the locally convex
topology generated by all amplitude seminorms.  For the Heisenberg
algebra $\cH_\kappa$, the sewing envelope is the symmetric algebra
of the Bergman space $A^2(D)$ of the disk
(Theorem~\ref{thm:heisenberg-sewing}\,(i)), and the genus-$g$
partition function is the Fredholm determinant of an explicit
trace-class sewing operator on $A^2(D)$.  The HS-sewing criterion
(Theorem~\ref{thm:general-hs-sewing}) extends the convergence to
the entire standard landscape of universal algebras.  On the proved
Gaussian scalar lane
\textup{(}Definition~\ref{def:scalar-lane}\textup{)},
the scalar package is a $\kappa$-multiple of the Heisenberg one, so
the Fredholm structure propagates there; beyond that lane, higher-arity
corrections need not be Fredholm.

The section is organized as follows.
\begin{enumerate}[label=\textup{(\arabic*)}]
\item Sewing envelopes and Bergman spaces:
  the analytic framework for one-particle reduction
  (\S\ref{subsec:thqg-X-sewing-bergman}).
\item The Heisenberg sewing theorem in full:
  four clauses with complete proofs
  (\S\ref{subsec:thqg-X-heisenberg-sewing}).
\item One-particle reduction, trace-class estimates,
  and the Dedekind eta derivation
  (\S\ref{subsec:thqg-X-one-particle}).
\item Extension to class-G algebras via scalar saturation
  (\S\ref{subsec:thqg-X-class-G}).
\item Non-Fredholm corrections for classes L, C, M
  (\S\ref{subsec:thqg-X-non-fredholm}).
\item Analytic aspects: the analytic bar coalgebra,
  coderived categories, and open problems
  (\S\ref{subsec:thqg-X-analytic-aspects}).
\end{enumerate}

% =====================================================================
%
%  SECTION 1: SEWING ENVELOPES AND BERGMAN SPACES
%
% =====================================================================

\subsection{Sewing envelopes and Bergman spaces}
\label{subsec:thqg-X-sewing-bergman}

\subsubsection{The Bergman space of a disk}

\begin{definition}[Bergman space of a disk]%
\label{def:thqg-X-bergman-space}%
\index{Bergman space|textbf}%
Let $D = \{z \in \C : |z| < 1\}$ be the open unit disk.
The \emph{Bergman space} $A^2(D)$ is the Hilbert space of
holomorphic functions that are square-integrable with respect
to the area measure:
\begin{equation}\label{eq:thqg-X-bergman-norm}
A^2(D)
\;:=\;
\Bigl\{f \in \cO(D) : \|f\|^2_{A^2}
= \int_D |f(z)|^2\,dA(z) < \infty\Bigr\},
\end{equation}
where $dA(z) = \frac{i}{2}\,dz \wedge d\bar{z}
= r\,dr\,d\theta$ in polar coordinates.  The monomials
$e_n(z) = \sqrt{(n+1)/\pi}\, z^n$ ($n \geq 0$)
form an orthonormal basis, and the reproducing kernel is
\begin{equation}\label{eq:thqg-X-bergman-kernel}
\Bergker_D(z,w)
\;=\;
\sum_{n=0}^{\infty} e_n(z)\overline{e_n(w)}
\;=\;
\frac{1}{\pi(1 - z\bar{w})^2}.
\end{equation}
\end{definition}

\begin{remark}[Conformal weight and sections of $K^{-1/2}$]%
\label{rem:thqg-X-bergman-conformal}%
The prime form $\primeform(z_1,z_2)$ on a Riemann surface is a section
of $K^{-1/2} \boxtimes K^{-1/2}$, \emph{not} $K^{+1/2}$.
The Bergman kernel
$\Bergker_D(z,w)$ is a section of $K \boxtimes \bar{K}$, consistent
with its conformal weight $(1,1)$.  The mode--Bergman correspondence
(Remark~\ref{rem:heisenberg-mode-bergman}) maps Fock modes to
Bergman monomials at weight $0$, absorbing the conformal weight
into the $\sqrt{n+1}$ prefactor.
\end{remark}

\begin{definition}[Bergman space with collar parameter]%
\label{def:thqg-X-bergman-collar}%
\index{Bergman space!collar parameter}%
For $0 < q < 1$, define the \emph{weighted Bergman space}
\begin{equation}\label{eq:thqg-X-bergman-weighted}
A^2_q(D)
\;:=\;
\Bigl\{f \in \cO(D) :
\|f\|^2_{q} := \sum_{n=0}^{\infty}
q^n \,|a_n|^2 \,(n+1)^{-1} < \infty
\;\text{where}\; f = \sum_{n=0}^{\infty} a_n z^n\Bigr\}.
\end{equation}
The collar parameter $q = e^{-t}$ ($t > 0$ the collar length)
encodes the sewing parameter for pair-of-pants decompositions.
The inclusion $A^2_{q'}(D) \hookrightarrow A^2_q(D)$ for $q' < q$
is trace class.
\end{definition}

\begin{lemma}[Bergman projection; \ClaimStatusProvedElsewhere]%
\label{lem:thqg-X-bergman-projection}%
\index{Bergman projection}%
The orthogonal projection
$P \colon L^2(D, dA) \to A^2(D)$ is the integral operator
\[
(Pf)(z)
\;=\;
\int_D \Bergker_D(z,w)\,f(w)\,dA(w)
\;=\;
\frac{1}{\pi}\int_D \frac{f(w)}{(1-z\bar{w})^2}\,dA(w).
\]
The projection satisfies $P^2 = P$, $P^* = P$, and
$\ker P = (\bar{\partial}\text{-exact})$.
\end{lemma}

\subsubsection{The sewing envelope}

\begin{definition}[Sewing envelope; review]%
\label{def:thqg-X-sewing-envelope-review}%
\index{sewing envelope!review}%
Recall from Definition~\ref{def:sewing-envelope} that
the \emph{sewing envelope} $\sewenv$ of an algebraic chiral
core $\algcore$ is the Hausdorff completion
for the locally convex topology generated by all amplitude
seminorms $p_{\Sigma,\xi,\eta}(a)
= |\langle \eta, A_\Sigma^{\mathrm{alg}}(\xi_1 \otimes \cdots
\otimes a \otimes \cdots \otimes \xi_m)\rangle|$,
where the supremum ranges over all finite conformally flat
bordisms $\Sigma$ with parametrized collars and all
test states $\xi, \eta$.  By Proposition~\ref{prop:sewing-universal-property},
$\sewenv$ has the universal property: every complete locally convex
realization factors uniquely through~$\sewenv$.
\end{definition}

\begin{proposition}[Sewing envelope: conformal weight filtration;
\ClaimStatusProvedHere]%
\label{prop:thqg-X-sewing-filtration}%
\index{sewing envelope!filtration}%
Let $\cA$ be a positive-energy chiral algebra with
conformal weight decomposition
$\algcore = \bigoplus_{n \geq 0} \algcore_n$, where
$\dim \algcore_n < \infty$ for all~$n$.  Then:
\begin{enumerate}[label=\textup{(\roman*)}]
\item The sewing envelope inherits a filtration
  $F^{\leq N}\sewenv
  = \overline{\bigoplus_{n \leq N} \algcore_n}$
  whose associated graded recovers the algebraic core:
  $\mathrm{gr}\,\sewenv \cong \algcore$.
\item For each $0 < q < 1$, the $q$-weighted completion
  $\sewenv_q := \widehat{\bigoplus}_{n \geq 0} q^{n/2}\,\algcore_n$
  is a Banach algebra for the operator norm of
  $A^2_q$-bounded amplitudes.
\item The sewing envelope is the projective limit
  $\sewenv = \varprojlim_{q \to 1^-} \sewenv_q$.
\end{enumerate}
\end{proposition}

\begin{proof}
Part~(i): The amplitude seminorms $p_{\Sigma,\xi,\eta}$ are
continuous with respect to the conformal weight grading because
each amplitude map preserves total conformal weight up to a shift
determined by the Euler characteristic of~$\Sigma$.
Passage to the associated graded undoes the completion, recovering
the algebraic core.

Part~(ii): The $q$-weighting assigns norm $q^{n/2}$ to elements
of weight~$n$.  The pair-of-pants amplitude
$m \colon \algcore_a \otimes \algcore_b \to \bigoplus_c \algcore_c$
satisfies $|m^c_{a,b}| \leq K(a+b+c+1)^N$ by polynomial OPE growth.
The $q$-weighted operator norm is
$\|m\|_q \leq K \sum_c q^{c/2}(a+b+c+1)^N < \infty$
for fixed $a,b$ and $q < 1$, so $\sewenv_q$ is submultiplicative.
Completeness follows from the completeness of the
$q$-weighted $\ell^2$ space.

Part~(iii): The sewing topology is generated by
$\{p_{\Sigma,\xi,\eta}\}$, which are bounded by $q$-weighted
norms for any $q$ sufficiently close to~$1$ (depending on the
collar length of~$\Sigma$).  Hence $\sewenv$ is the intersection
$\bigcap_{q < 1} \sewenv_q = \varprojlim_{q} \sewenv_q$.
\end{proof}

\subsubsection{The Heisenberg sewing envelope}

\begin{proposition}[Heisenberg sewing envelope = $\operatorname{Sym} A^2(D)$;
\ClaimStatusProvedElsewhere{} Moriwaki~\cite{Moriwaki26b}]%
\label{prop:thqg-X-heisenberg-sewing-envelope}%
\index{Heisenberg!sewing envelope}%
\index{sewing envelope!Heisenberg}%
The sewing envelope of the algebraic Heisenberg VOA
$\cH_\kappa$ with level~$\kappa$ is
\begin{equation}\label{eq:thqg-X-heisenberg-sew}
\cH_\kappa^{\mathrm{sew}}
\;\cong\;
\operatorname{Sym} A^2(D),
\end{equation}
the symmetric (bosonic) Fock space over the Bergman space of
the disk.  The isomorphism is given by the mode--Bergman
correspondence
\begin{equation}\label{eq:thqg-X-mode-bergman-map}
\Theta(a_{-n-1}\mathbf{1})
\;=\;
\sqrt{(n+1)\kappa}\; z^n
\;\in\; A^2(D),
\end{equation}
extended multiplicatively by symmetrization:
$\Theta(a_{-n_1-1} \cdots a_{-n_k-1}\mathbf{1})
= \operatorname{Sym}(\Theta(a_{-n_1-1}\mathbf{1})
\otimes \cdots \otimes
\Theta(a_{-n_k-1}\mathbf{1}))$.
\end{proposition}

\begin{proof}
This is the ind-Hilbert identification of
Moriwaki~\cite{Moriwaki26b}, which constructs a conformally
flat 2-disk algebra structure on $\operatorname{Sym} A^2(D)$
and verifies that it is isomorphic to the Hilbert completion
of the Heisenberg Fock space for the sewing-amplitude topology.
The $\sqrt{(n+1)\kappa}$ normalization ensures that
$\|\Theta(a_{-n-1}\mathbf{1})\|^2 = \kappa$, matching the
level of the two-point function
$\langle a(z)\,a(w)\rangle = \kappa/(z-w)$.
\end{proof}

\begin{remark}[Tensor product structure]%
\label{rem:thqg-X-tensor-sewing}%
The Fock space $\operatorname{Sym} A^2(D)$ has the standard
$n$-particle decomposition
\[
\operatorname{Sym} A^2(D)
\;=\;
\bigoplus_{n=0}^{\infty}
\operatorname{Sym}^n A^2(D),
\]
where $\operatorname{Sym}^0 = \C$ (vacuum sector) and
$\operatorname{Sym}^n$ is the $n$-th symmetric power.
The dimension of the weight-$N$ subspace in
$\operatorname{Sym} A^2(D)$ is the number of partitions
$p(N)$, matching the Heisenberg character
$\chi_{\cH}(q) = \prod_{n=1}^{\infty}(1-q^n)^{-1}
= q^{-1/24}/\dedeta(q)$.
\end{remark}

\subsubsection{Sewing operators on Bergman spaces}

\begin{definition}[Sewing operator]%
\label{def:thqg-X-sewing-operator}%
\index{sewing operator|textbf}%
\index{Bergman space!sewing operator}%
Let $\Sigma$ be a compact Riemann surface obtained by sewing
two disks $D_1, D_2$ along their boundary circles via the
identification $z_1 z_2 = q$ for $0 < q < 1$.  The
\emph{sewing operator} $\sewop_q \colon A^2(D) \to A^2(D)$ is
defined on the orthonormal basis by
\begin{equation}\label{eq:thqg-X-sewing-eigenvalues}
\sewop_q\, e_n \;=\; q^{n+1}\, e_n,
\qquad n \geq 0.
\end{equation}
Equivalently, $\sewop_q$ acts by multiplication by
$|z|^2 \cdot q$ on the Bergman kernel:
\begin{equation}\label{eq:thqg-X-sewing-kernel}
(\sewop_q f)(z)
\;=\;
\int_D \sewop_q(z,w)\,f(w)\,dA(w),
\qquad
\sewop_q(z,w)
\;=\;
\sum_{n=0}^{\infty}
q^{n+1}\,e_n(z)\overline{e_n(w)}.
\end{equation}
\end{definition}

\begin{lemma}[Schatten class properties of $\sewop_q$;
\ClaimStatusProvedHere]%
\label{lem:thqg-X-sewing-schatten}%
\index{sewing operator!Schatten class}%
For $0 < q < 1$, the sewing operator $\sewop_q$ on $A^2(D)$
satisfies:
\begin{enumerate}[label=\textup{(\roman*)}]
\item \emph{Trace class.}
  $\sewop_q \in \trcl(A^2(D))$ with
  \begin{equation}\label{eq:thqg-X-trace-norm}
  \|\sewop_q\|_1
  \;=\;
  \sum_{n=0}^{\infty} q^{n+1}
  \;=\;
  \frac{q}{1-q}.
  \end{equation}
\item \emph{Hilbert--Schmidt.}
  $\sewop_q \in \hscl(A^2(D))$ with
  \begin{equation}\label{eq:thqg-X-hs-norm}
  \|\sewop_q\|_{\hscl}^2
  \;=\;
  \sum_{n=0}^{\infty} q^{2(n+1)}
  \;=\;
  \frac{q^2}{1-q^2}.
  \end{equation}
\item \emph{$p$-Schatten.}
  For all $p \geq 1$,
  $\sewop_q \in \mathcal{L}^p(A^2(D))$ with
  $\|\sewop_q\|_p^p = q^p/(1-q^p)$.
\item \emph{Spectral gap.}
  $\|\sewop_q\| = q$ (operator norm), and the spectral
  radius satisfies $r(\sewop_q) = q < 1$.
\end{enumerate}
\end{lemma}

\begin{proof}
The eigenvalues of $\sewop_q$ are $\lambda_n = q^{n+1}$ for
$n \geq 0$.  Since $0 < q < 1$, the eigenvalues are positive
and summable:
\[
\|\sewop_q\|_p^p
= \sum_{n=0}^{\infty} |\lambda_n|^p
= \sum_{n=0}^{\infty} q^{p(n+1)}
= \frac{q^p}{1-q^p}.
\]
For $p = 1$ (trace norm): $q/(1-q)$.
For $p = 2$ (HS norm squared): $q^2/(1-q^2)$.
The operator norm is $\|\sewop_q\| = \sup_n q^{n+1} = q$.
\end{proof}

\begin{definition}[Genus-$g$ sewing kernel]%
\label{def:thqg-X-genus-g-sewing}%
\index{sewing operator!genus $g$}%
A genus-$g$ surface $\Sigma_g$ with $n$ punctures admits a
pair-of-pants decomposition into $2g - 2 + n$ pairs of pants,
with $3g - 3 + n$ sewing circles.  Each sewing circle
$C_\alpha$ ($\alpha = 1, \ldots, 3g-3+n$) contributes a sewing
parameter $q_\alpha = e^{-t_\alpha + i\theta_\alpha}$ and
a sewing operator $\sewop_{q_\alpha}$ on $A^2(D)$.
The \emph{genus-$g$ sewing kernel} is the composite operator
\begin{equation}\label{eq:thqg-X-genus-g-kernel}
\sewker_g
\;=\;
\prod_{\alpha=1}^{3g-3}
\sewop_{q_\alpha}
\;\in\;
\trcl(A^2(D)^{\otimes g}),
\end{equation}
where the product is over the $3g-3$ internal sewing circles
of a genus-$g$ surface without punctures.
More precisely, $\sewker_g$ acts on the tensor product of
$g$ copies of $A^2(D)$ (one for each handle),
with each $\sewop_{q_\alpha}$ acting on the appropriate
tensor factor.
\end{definition}

\begin{remark}[Dependence on pants decomposition]%
\label{rem:thqg-X-pants-independence}%
\index{pants decomposition!independence}%
The Fredholm determinant $\det(1 - \sewker_g)$ is independent
of the choice of pair-of-pants decomposition, although the
intermediate operator $\sewker_g$ depends on the decomposition.
This independence follows from the sewing axiom of the
chain-level modular-functor package
(Theorem~\ref{thm:chain-modular-functor}\,(iv)):
different decompositions are related by chain homotopies
that preserve the trace.
\end{remark}

\subsubsection{Composition operators from the prime form}

\begin{definition}[Bergman composition operator]%
\label{def:thqg-X-bergman-composition}%
\index{composition operator!Bergman}%
Let $\Sigma_g$ be a genus-$g$ Riemann surface with period
matrix $\Omega \in \Siegel_g$ (the Siegel upper half-space),
and let $D_j \subset \Sigma_g$ be small coordinate disks around
marked points $p_j$.  The \emph{Bergman composition operator}
$R \colon A^2(D_{\mathrm{out}}) \to
A^2(D_{\mathrm{in},1}) \otimes A^2(D_{\mathrm{in},2})$
for a pair-of-pants $P$ with one outgoing and two incoming
boundaries is
\begin{equation}\label{eq:thqg-X-bergman-composition}
R_{n,m}
\;=\;
\frac{1}{(2\pi i)^2}
\oint_{|z|=1} \oint_{|w|=1}
z^{-n-1}\,w^{-m-1}\,
G_P(z,w)\,dz\,dw,
\end{equation}
where $G_P(z,w) = \partial_z \log \primeform(z,w;\Omega)$ is the
chiral Green function on the pair of pants.
\end{definition}

\begin{lemma}[Exponential decay of composition coefficients;
\ClaimStatusProvedHere]%
\label{lem:thqg-X-composition-decay}%
\index{composition operator!exponential decay}%
The matrix coefficients of the Bergman composition operator satisfy
\begin{equation}\label{eq:thqg-X-composition-bound}
|R_{n,m}|
\;\leq\;
C\,q^{(n+m)/2}
\end{equation}
for a constant $C$ depending only on the collar length
$t = -\log q$ of the pair of pants.  In particular,
$R$ is Hilbert--Schmidt with
$\|R\|_{\hscl}^2 \leq C^2/(1-q)^2$.
\end{lemma}

\begin{proof}
The Green function $G_P(z,w)$ on the pair of pants is holomorphic
for $|z| < 1$, $|w| < 1$ away from the diagonal $z = w$.
By the collar theorem for hyperbolic surfaces, there exists a
collar of width $t/2$ around each boundary circle, so $G_P$ extends
holomorphically to $|z| < e^{t/2}$, $|w| < e^{t/2}$.
The Cauchy estimate on the annulus gives
$|R_{n,m}| \leq \sup_{|z|=e^{t/4}} |G_P|\, e^{-t(n+m)/4}
= C\,q^{(n+m)/4}$ with $q = e^{-t}$.
The sharper bound $q^{(n+m)/2}$ follows by moving the contour
further into the collar.

For the HS norm:
\[
\|R\|_{\hscl}^2
= \sum_{n,m \geq 0} |R_{n,m}|^2
\leq C^2 \sum_{n,m \geq 0} q^{n+m}
= \frac{C^2}{(1-q)^2}.
\qedhere
\]
\end{proof}

\begin{proposition}[Second quantization;
\ClaimStatusProvedHere]%
\label{prop:thqg-X-second-quantization}%
\index{second quantization!Bergman}%
Let $T \colon A^2(D) \to A^2(D)$ be a trace-class operator
with $\|T\| < 1$.  The \emph{second quantization}
$\secquant(T) \colon \operatorname{Sym} A^2(D) \to
\operatorname{Sym} A^2(D)$ is defined by
\begin{equation}\label{eq:thqg-X-second-quant}
\secquant(T)|_{\operatorname{Sym}^n A^2(D)}
\;=\;
T^{\otimes n}\big|_{\operatorname{Sym}^n},
\end{equation}
with $\secquant(T)|_{\operatorname{Sym}^0} = \id$.
Then:
\begin{enumerate}[label=\textup{(\roman*)}]
\item $\secquant(T)$ is trace class on $\operatorname{Sym} A^2(D)$
  with
  \begin{equation}\label{eq:thqg-X-second-quant-trace}
  \operatorname{Tr}_{\operatorname{Sym}}(\secquant(T))
  \;=\;
  \det(1 - T)^{-1}.
  \end{equation}
\item More generally, for the Fredholm determinant:
  \begin{equation}\label{eq:thqg-X-fredholm-second-quant}
  \det_{\operatorname{Sym}}(1 - \secquant(T))
  \;=\;
  \exp\!\Bigl(-\sum_{k=1}^{\infty}
  \frac{1}{k}\operatorname{Tr}(T^k)\Bigr)
  \;=\;
  \det(1 - T).
  \end{equation}
\end{enumerate}
\end{proposition}

\begin{proof}
Part~(i): The trace on $\operatorname{Sym}^n A^2(D)$ is computed
using an orthonormal basis of symmetric tensors.
For eigenvalues $\lambda_j$ of~$T$:
\[
\operatorname{Tr}_{\operatorname{Sym}^n}(T^{\otimes n})
= \sum_{j_1 \leq \cdots \leq j_n}
\lambda_{j_1} \cdots \lambda_{j_n}
= h_n(\lambda_1, \lambda_2, \ldots),
\]
where $h_n$ is the complete symmetric polynomial.
Summing over $n \geq 0$:
\[
\sum_{n=0}^{\infty}
\operatorname{Tr}_{\operatorname{Sym}^n}(T^{\otimes n})
= \sum_{n=0}^{\infty} h_n(\lambda)
= \prod_{j} \frac{1}{1 - \lambda_j}
= \det(1 - T)^{-1}.
\]
The series converges absolutely because $T$ is trace class
and $\|T\| < 1$, so $\sum_j |\lambda_j| < \infty$ and
$|\lambda_j| < 1$ for all~$j$.

Part~(ii): $\log\det(1-T) = \sum_j \log(1-\lambda_j)
= -\sum_{k=1}^{\infty} \frac{1}{k}\sum_j \lambda_j^k
= -\sum_{k=1}^{\infty} \frac{1}{k}\operatorname{Tr}(T^k)$.
\end{proof}

% =====================================================================
%
%  SECTION 2: THE HEISENBERG SEWING THEOREM
%
% =====================================================================

\subsection{The Heisenberg sewing theorem: full development}
\label{subsec:thqg-X-heisenberg-sewing}

\subsubsection{Statement and architecture}

The Heisenberg sewing theorem
(Theorem~\ref{thm:heisenberg-sewing})
has four clauses.  We now develop the complete proofs, starting
from the mode--Bergman correspondence and building to the
Fredholm determinant formula for genus-$g$ partition functions.

\begin{theorem}[Heisenberg sewing theorem: full development;
\ClaimStatusProvedHere]%
\label{thm:thqg-X-heisenberg-sewing-full}%
\index{Heisenberg!sewing theorem!full development}%
Let $\cH_\kappa$ be the Heisenberg algebra with level $\kappa > 0$,
equipped with the mode--Bergman correspondence~$\Theta$
\textup{(}Remark~\textup{\ref{rem:heisenberg-mode-bergman}}\textup{)}.
Then:
\begin{enumerate}[label=\textup{(\Roman*)}]
\item \textbf{Sewing envelope identification.}
  The sewing envelope of $\cH_\kappa$ is
  $\cH_\kappa^{\mathrm{sew}} \cong
  \operatorname{Sym} A^2(D)$.
\item \textbf{Bar differential closure.}
  The algebraic bar differential extends by closure to the
  Gaussian collision operator on Bergman tensors:
  \begin{equation}\label{eq:thqg-X-bar-closure}
  d_{\mathrm{bar}}^{\mathrm{an}}
  \;=\;
  \overline{d_{\mathrm{bar}}^{\mathrm{alg}}}
  ^{\,\graphnorm}
  \;=\;
  \sum_{i < j}
  \kappa\;\operatorname{Res}_{z_i = z_j}
  \frac{1}{(z_i - z_j)}\,
  \partial_{z_i z_j}.
  \end{equation}
\item \textbf{Pair-of-pants sewing = second quantization.}
  The pair-of-pants amplitude on
  $\operatorname{Sym} A^2(D)$ is the second quantization
  $\secquant(R)$ of the one-particle Bergman restriction
  operator
  $R \colon A^2(D_{\mathrm{out}}) \to
  A^2(D_{\mathrm{in},1}) \otimes A^2(D_{\mathrm{in},2})$.
\item \textbf{Fredholm determinant formula.}
  Every genus-$g$ closed amplitude is a Fredholm determinant:
  the genus-$g$ partition function of $\cH_\kappa$ is
  \begin{equation}\label{eq:thqg-X-fredholm-formula}
  Z_g(\cH_\kappa;\,\Omega)
  \;=\;
  \bigl(\det\operatorname{Im}\Omega\bigr)^{-\kappa/2}
  \cdot
  \bigl(\zetareg\Delta_{\Sigma_g}\bigr)^{-\kappa/2},
  \end{equation}
  where $\Omega$ is the period matrix of $\Sigma_g$ and
  $\zetareg\Delta$ is the zeta-regularized determinant of the
  Laplacian.
\end{enumerate}
\end{theorem}

The proof occupies the remainder of this subsection.

\subsubsection{Clause I: sewing envelope identification}

\begin{proof}[Proof of Clause~I]
The mode--Bergman map $\Theta$ of
Remark~\ref{rem:heisenberg-mode-bergman} defines a dense
algebraic embedding
$\Theta \colon \Fock_\kappa^{\mathrm{fin}} \hookrightarrow
\operatorname{Sym} A^2(D)$
from the algebraic Fock space (finite linear combinations of
mode monomials) into the symmetric algebra of the Bergman space.

\emph{Step 1: algebraic amplitudes.}
On the algebraic Fock core, every amplitude
$A_\Sigma^{\mathrm{alg}}$ is a finite sum of Wick contractions.
By Wick's theorem for the free boson, the amplitude of
$a_{-n_1-1}\cdots a_{-n_k-1}\mathbf{1}$ on a surface~$\Sigma$
with collar coordinates is a polynomial in the mode
coefficients of the propagator $G_\Sigma$, which by the
mode--Bergman dictionary become matrix coefficients of the
Bergman composition operator $R$.

\emph{Step 2: continuity for the sewing topology.}
The sewing seminorms $p_{\Sigma,\xi,\eta}$ on $\Fock_\kappa^{\mathrm{fin}}$
are controlled by the Bergman norms: for
$v = \Theta(a_{-n_1-1}\cdots a_{-n_k-1}\mathbf{1})
\in \operatorname{Sym}^k A^2(D)$,
\begin{equation}\label{eq:thqg-X-sewing-seminorm-bound}
p_{\Sigma,\xi,\eta}(v)
\;\leq\;
\|\xi\| \cdot \|\eta\| \cdot \|v\|_{\operatorname{Sym}^k}
\cdot \prod_\alpha \|\sewop_{q_\alpha}\|,
\end{equation}
where the product is over sewing circles of~$\Sigma$.
Since each $\sewop_{q_\alpha}$ is bounded with
$\|\sewop_{q_\alpha}\| = q_\alpha < 1$, the sewing topology is
weaker than the $\operatorname{Sym} A^2(D)$ topology.

\emph{Step 3: completeness.}
Conversely, any Cauchy sequence in the sewing topology is Cauchy
in $\operatorname{Sym} A^2(D)$ (because the latter norm
controls each sewing seminorm).  Hence the completions coincide:
$\cH_\kappa^{\mathrm{sew}} \cong \operatorname{Sym} A^2(D)$.
The full argument is the identification of Moriwaki~\cite{Moriwaki26b},
which constructs the conformally flat 2-disk algebra structure
on $\operatorname{Sym} A^2(D)$ and verifies the universal
property of the sewing envelope.
\end{proof}

\subsubsection{Clause II: bar differential closure}

\begin{proof}[Proof of Clause~II]
The algebraic bar differential on $\barB_n(\cH_\kappa)$ is
\[
d_{\mathrm{bar}}^{\mathrm{alg}}
= \sum_{i < j}
\operatorname{Res}_{z_i = z_j}
\circ\, Y_{ij},
\]
where $Y_{ij}$ is the chiral multiplication
$Y(a_{-n-1}\mathbf{1}, z_i - z_j)
= \sum_m a_m (z_i - z_j)^{-m-1}$.
For the Heisenberg algebra, the OPE is
$a(z_i)\,a(z_j) \sim \kappa/(z_i - z_j)$,
so $Y_{ij}$ acts on the one-particle Bergman space as
multiplication by $\kappa/(z_i - z_j)$ followed by the residue
at $z_i = z_j$.

\emph{Step 1: densely defined operator.}
On $\operatorname{Sym}^n A^2(D)$, the operator
$d_{\mathrm{bar}}^{\mathrm{alg}}$ is densely defined on the
algebraic Fock core.  Its action on monomials is
\begin{align}
d_{\mathrm{bar}}(z^{n_1} \cdots z^{n_k})
&= \kappa \sum_{i < j}
\operatorname{Res}_{z_i = z_j}
\frac{z_1^{n_1} \cdots z_k^{n_k}}{z_i - z_j}
\notag\\
&= \kappa \sum_{i < j}
\sum_{p=0}^{n_j-1}
z_j^{n_i + n_j - 1 - p}\, z_j^p
\cdot \prod_{\ell \neq i,j} z_\ell^{n_\ell},
\label{eq:thqg-X-bar-monomial-action}
\end{align}
where the residue replaces $z_i$ by $z_j$ with the standard
contraction.

\emph{Step 2: closability.}
The operator $d_{\mathrm{bar}}^{\mathrm{alg}}$ is closable
on $\operatorname{Sym}^n A^2(D)$.  This follows from the
fact that the formal adjoint $(d_{\mathrm{bar}}^{\mathrm{alg}})^*$
is also densely defined (it is the creation operator that reverses
the contraction), so the closure exists by von Neumann's theorem.

\emph{Step 3: identification of the closure.}
The closure is the Gaussian collision operator: it acts on
$f_1 \otimes \cdots \otimes f_n \in \operatorname{Sym}^n A^2(D)$
by
\[
\overline{d}_{\mathrm{bar}}(f_1 \otimes \cdots \otimes f_n)
= \kappa \sum_{i < j}
\operatorname{Res}_{z_i = z_j}
\frac{f_i(z_i)\,f_j(z_j)}{z_i - z_j}
\cdot \prod_{\ell \neq i,j} f_\ell(z_\ell),
\]
where the residue is computed by the Bergman projection
of $f_i(z)\,f_j(z) / (z_i - z_j)$ restricted to the diagonal.
The graph norm of this operator is bounded by the Bergman
norm of the input, establishing $\overline{d}_{\mathrm{bar}}$
as the unique closed extension.
\end{proof}

\subsubsection{Clause III: pair-of-pants = second quantization}

\begin{proof}[Proof of Clause~III]
Consider a pair of pants $P$ with one outgoing boundary $\partial_{\mathrm{out}}$
and two incoming boundaries $\partial_{\mathrm{in},1}$, $\partial_{\mathrm{in},2}$.
Each boundary is parametrized by a circle, and the sewing
parameters $(q_1, q_2)$ are the collar parameters.

\emph{Step 1: one-particle amplitude.}
On the one-particle space $A^2(D)$, the pair-of-pants amplitude
is the Bergman composition operator
$R \colon A^2(D_{\mathrm{out}}) \to
A^2(D_{\mathrm{in},1}) \otimes A^2(D_{\mathrm{in},2})$
of Definition~\ref{def:thqg-X-bergman-composition}.
Its matrix coefficients $R_{n,m}$ are the contour integrals
of the pair-of-pants Green function $G_P(z,w)$.

\emph{Step 2: Wick's theorem.}
On the full Fock space $\operatorname{Sym} A^2(D)$, the
Heisenberg amplitudes are determined by Wick's theorem:
every $n$-point function is a sum over pairings of two-point
functions.  In the Bergman picture, the $n$-particle amplitude
on $P$ is
\begin{equation}\label{eq:thqg-X-wick-bergman}
A_P^{(n)}(f_1 \otimes \cdots \otimes f_n)
= \sum_{\sigma \in S_n}
\prod_{j=1}^{n}
\langle f_j, R\,g_{\sigma(j)}\rangle,
\end{equation}
where $g_{\sigma(j)}$ are the incoming modes.  This is precisely
the formula for the second quantization:
\[
A_P^{(n)}
= \secquant(R)\big|_{\operatorname{Sym}^n}.
\]

\emph{Step 3: total amplitude.}
Summing over particle number gives the full pair-of-pants
amplitude:
\begin{equation}\label{eq:thqg-X-pants-second-quant}
A_P
= \bigoplus_{n=0}^{\infty}
A_P^{(n)}
= \secquant(R)
\colon \operatorname{Sym} A^2(D_{\mathrm{out}})
\to
\operatorname{Sym}(A^2(D_{\mathrm{in},1})
\oplus A^2(D_{\mathrm{in},2})).
\end{equation}
The convergence of this series follows from the exponential
decay $|R_{n,m}| \leq C q^{(n+m)/2}$
(Lemma~\ref{lem:thqg-X-composition-decay}).
\end{proof}

\subsubsection{Clause IV: Fredholm determinant}

\begin{proof}[Proof of Clause~IV]
A genus-$g$ surface $\Sigma_g$ without punctures is obtained by
sewing together $2g-2$ pairs of pants along $3g-3$ sewing circles.
Each sewing circle contributes a trace over the one-particle
Hilbert space, and by Wick's theorem the total amplitude factors
through the one-particle operators.

\emph{Step 1: genus-$g$ amplitude as iterated trace.}
The genus-$g$ partition function is
\begin{equation}\label{eq:thqg-X-genus-g-trace}
Z_g(\cH_\kappa; \Omega)
= \operatorname{Tr}_{\operatorname{Sym} A^2}
\Bigl(\prod_{\alpha=1}^{3g-3}
\secquant(\sewop_{q_\alpha})^{\kappa}\Bigr),
\end{equation}
where the product is over all sewing circles, $\sewop_{q_\alpha}$
is the one-particle sewing operator, and the exponent $\kappa$
reflects the level (one Fock space per unit of central charge).

\emph{Step 2: trace factorization.}
By Proposition~\ref{prop:thqg-X-second-quantization}\,(i):
\[
\operatorname{Tr}_{\operatorname{Sym}}(\secquant(T))
= \det(1 - T)^{-1}
\]
for any trace-class $T$ with $\|T\| < 1$.  Applying this to each
sewing circle:
\begin{equation}\label{eq:thqg-X-genus-g-det}
Z_g(\cH_\kappa; \Omega)
= \prod_{\alpha=1}^{3g-3}
\det\bigl(1 - \sewop_{q_\alpha}\bigr)^{-\kappa}.
\end{equation}

\emph{Step 3: from pants parameters to period matrix.}
The sewing parameters $\{q_\alpha\}$ are the Fenchel--Nielsen
coordinates of $\Sigma_g$; they determine and are determined by
the period matrix $\Omega \in \Siegel_g$.  The Polyakov--Alvarez
formula \cite{BK86} relates the product of one-particle
determinants to the period matrix and the zeta-regularized
determinant of the Laplacian:
\begin{equation}\label{eq:thqg-X-polyakov}
\prod_{\alpha=1}^{3g-3}
\det(1 - \sewop_{q_\alpha})^{-1}
= (\det\operatorname{Im}\Omega)^{-1/2}
\cdot (\zetareg\Delta_{\Sigma_g})^{-1/2}.
\end{equation}
This is the holomorphic factorization of the string measure
of Belavin--Knizhnik~\cite{BK86}: the modular-invariant
expression factors into a product of a holomorphic and an
anti-holomorphic part, and the holomorphic part is
$\det(1 - \sewker_g)^{-1}$ up to the period-matrix prefactor.

\emph{Step 4: level $\kappa$.}
Raising to the power $\kappa$ (i.e., tensoring $\kappa$ copies
of the one-particle Bergman space, or equivalently working at
level $\kappa$):
\begin{equation}\label{eq:thqg-X-final-fredholm}
Z_g(\cH_\kappa; \Omega)
= (\det\operatorname{Im}\Omega)^{-\kappa/2}
\cdot (\zetareg\Delta_{\Sigma_g})^{-\kappa/2}.
\end{equation}
This completes the proof of Clause~IV and of the full theorem.
\end{proof}

\begin{remark}[Functional determinant--Fredholm bridge]%
\label{rem:functional-det-fredholm-bridge}%
\index{functional determinant!Fredholm bridge}%
Polyakov's functional determinant $\zetareg\Delta$ is defined via
zeta regularization: the formal product $\prod_n \lambda_n$ of
eigenvalues of the Laplacian is regularized by analytic continuation
of $\zeta_\Delta(s) = \sum_n \lambda_n^{-s}$ to $s = 0$.
The Heisenberg sewing theorem
(Theorem~\ref{thm:thqg-X-heisenberg-sewing-full})
proves that each genus-$g$ closed amplitude is a Fredholm determinant
$\det(1 - \sewker_g)^{-\kappa}$ on the one-particle Bergman space
$A^2(D)$, where $\sewker_g$ is the composition operator encoding
pair-of-pants gluing.  The Hilbert--Schmidt condition on
the sewing operators ensures $\sewker_g$ is trace class, so the
Fredholm determinant converges absolutely.  This makes Polyakov's
formal manipulations rigorous: the zeta-regularized Laplacian
determinant \emph{is} the Fredholm determinant on the Bergman
one-particle space, via the identity
\eqref{eq:thqg-X-polyakov}.
\end{remark}

\begin{proposition}[Polyakov--Alvarez as shadow metric
specialization; \ClaimStatusProvedHere]%
\label{prop:polyakov-alvarez-shadow-specialization}%
\index{Polyakov--Alvarez formula!shadow metric specialization}%
The Polyakov--Alvarez formula at genus~$g$, which converts the
product of one-particle Fredholm determinants to period-matrix
data, is the specialization of the shadow metric
$Q_L(t)$ \textup{(}Definition~\textup{\ref{def:shadow-metric}}\textup{)}
to the Heisenberg line $L = \cH_\kappa$.
On this line, $Q_{\cH_\kappa}(t) = (2\kappa)^2$ \textup{(}constant,
since $\Delta = 0$ for the Gaussian class\textup{)}.
The Polyakov--Alvarez identity
\begin{equation}\label{eq:thqg-X-polyakov-shadow}
\prod_{n=1}^{\infty}
\det(1 - \sewker_g^{(n)})^{-\kappa}
\;=\;
\bigl(\det\operatorname{Im}\Omega_g\bigr)^{-\kappa/2}
\cdot
\bigl(\zetareg\Delta_{\Sigma_g}\bigr)^{-\kappa/2}
\end{equation}
expresses the factorization of the Fredholm determinant through
the period matrix $\Omega_g$, with the shadow metric providing the
proportionality constant $\kappa$.
\end{proposition}

\begin{proof}
The Heisenberg algebra $\cH_\kappa$ lies in the Gaussian class~$G$
(Definition~\ref{def:shadow-depth-classification}), so its shadow
tower terminates at arity~$2$: the cubic shadow $C = 0$, the quartic
contact invariant $Q^{\mathrm{contact}} = 0$, and the critical
discriminant $\Delta = 8\kappa S_4 = 0$.  Hence the shadow metric is
$Q_{\cH_\kappa}(t) = (2\kappa + 3\alpha t)^2 + 2\Delta\,t^2
= (2\kappa)^2$, a perfect square independent of $t$.
The genus-$g$ partition function is therefore determined entirely
by the arity-$2$ data $\kappa$: the Fredholm determinant
$\det(1 - \sewker_g)^{-\kappa}$ factors through the period matrix
via the classical Polyakov--Alvarez formula
\eqref{eq:thqg-X-polyakov}.  The proportionality constant is
$\kappa$ because each tensor factor of the Fock space
$\operatorname{Sym} A^2(D)$ contributes one copy of the
one-particle determinant $\det(1 - \sewker_g)^{-1}$.
\end{proof}

\begin{remark}[UV finiteness from compactness]%
\label{rem:uv-finiteness-compactness}%
\index{UV finiteness!compactness}%
No UV regularization is needed for the shadow obstruction tower genus expansion.
The compactness of $\Sigma_g$ guarantees: eigenvalues $\lambda_n > 0$
for $n \geq 1$, spectral zeta function $\zeta_\Delta(s)$
convergent for $\operatorname{Re}(s) > 1$ by the Weyl law
$\lambda_n \sim C n$ ($\dim = 2$), and smooth analytic continuation
to $s = 0$.  The Fredholm determinant $\det(1 - \sewker_g)$
converges because $\sewker_g$ is trace class, which follows from
the HS-sewing criterion
(Theorem~\ref{thm:thqg-X-heisenberg-sewing-full}).
This is the analytic reflection of the algebraic completeness of the
bar complex on compact Fulton--MacPherson spaces: the FM
compactification $\overline{\FM}_n(X)$ resolves all collision
singularities, so no UV subtraction is required at any stage of the
shadow obstruction tower.
\end{remark}

% =====================================================================
%
%  SECTION 3: ONE-PARTICLE REDUCTION
%
% =====================================================================

\subsection{One-particle reduction: trace-class operators
and the Dedekind eta}
\label{subsec:thqg-X-one-particle}

The one-particle reduction principle reduces the computation
of partition functions on the bosonic Fock space to Fredholm
determinants on the one-particle Bergman space.  This subsection
carries out the explicit computations at genus $1$, $2$, and $3$.

\subsubsection{Genus 1: the torus and the Dedekind eta}

\begin{theorem}[Genus-1 Fredholm determinant;
\ClaimStatusProvedHere]%
\label{thm:thqg-X-genus-1-fredholm}%
\index{Fredholm determinant!genus 1}%
\index{Dedekind eta function!from Fredholm determinant}%
The genus-$1$ partition function of $\cH_\kappa$ is
\begin{equation}\label{eq:thqg-X-genus-1-det}
Z_1(\cH_\kappa;\,\tau)
\;=\;
(\operatorname{Im}\tau)^{-\kappa/2}\,
\bigl|\dedeta(\tau)\bigr|^{-2\kappa},
\end{equation}
where the holomorphic part is
\begin{equation}\label{eq:thqg-X-genus-1-hol}
Z_1^{\mathrm{hol}}(\cH_\kappa;\,\tau)
\;=\;
\dedeta(\tau)^{-\kappa}
\;=\;
q^{-\kappa/24}\,
\prod_{n=1}^{\infty}
(1 - q^n)^{-\kappa}.
\end{equation}
\end{theorem}

\begin{proof}
A genus-$1$ surface $E_\tau = \C/(\Z + \tau\Z)$ is obtained
by sewing a single cylinder $\C^* \supset
\{q < |z| < 1\}$ with $q = e^{2\pi i \tau}$ ($\operatorname{Im}\tau > 0$).
The sewing circle is a single boundary component, and the
sewing operator is $\sewop_q$ on $A^2(D)$ with eigenvalues
$q^{n+1}$.

\emph{Step 1: Fock space trace.}
The partition function is the Fock-space trace:
\begin{align}
Z_1(\cH_1;\,\tau)
&= \operatorname{Tr}_{\operatorname{Sym} A^2(D)}
\bigl(\secquant(\sewop_q)\bigr)
\notag\\
&= \det(1 - \sewop_q)^{-1}
\quad\text{(Proposition~\ref{prop:thqg-X-second-quantization}\,(i))}
\notag\\
&= \prod_{n=0}^{\infty}
(1 - q^{n+1})^{-1}
\notag\\
&= \prod_{n=1}^{\infty}
(1 - q^n)^{-1}.
\label{eq:thqg-X-genus-1-product}
\end{align}

\emph{Step 2: Dedekind eta identification.}
The Dedekind eta function is
$\dedeta(\tau) = q^{1/24}\prod_{n=1}^{\infty}(1-q^n)$
(Definition~\ref{def:eta-function}).  Therefore:
\begin{equation}\label{eq:thqg-X-eta-identification}
\prod_{n=1}^{\infty}(1-q^n)^{-1}
= q^{1/24} / \dedeta(\tau)
= q^{1/24}\,\dedeta(\tau)^{-1}.
\end{equation}
Including the vacuum energy $q^{-c/24} = q^{-1/24}$ for
$c = 1$ (one free boson with central charge $1$ per unit of
$\kappa$) and the non-holomorphic prefactor
$(\operatorname{Im}\tau)^{-1/2}$ from the zero-mode
integral:
\[
Z_1(\cH_1;\,\tau)
= (\operatorname{Im}\tau)^{-1/2}\,
q^{-1/24}\,\bar{q}^{-1/24}\,
\prod_{n=1}^{\infty}
\frac{1}{|1-q^n|^2}
= (\operatorname{Im}\tau)^{-1/2}\,
|\dedeta(\tau)|^{-2}.
\]

\emph{Step 3: level $\kappa$.}
At level $\kappa$, the Fock space is the $\kappa$-fold tensor
product of the $\kappa = 1$ Fock space (or equivalently, the
sewing operator eigenvalues are raised to the power $\kappa$):
\[
Z_1(\cH_\kappa;\,\tau)
= (\operatorname{Im}\tau)^{-\kappa/2}\,
|\dedeta(\tau)|^{-2\kappa}.
\qedhere
\]
\end{proof}

\begin{remark}[Modular transformation]%
\label{rem:thqg-X-genus-1-modular}%
Under $\tau \mapsto -1/\tau$:
$\dedeta(-1/\tau) = \sqrt{-i\tau}\,\dedeta(\tau)$, so
\[
Z_1(\cH_\kappa;\,-1/\tau)
= |\tau|^{-\kappa}\,
(\operatorname{Im}\tau/|\tau|^2)^{-\kappa/2}\,
|\dedeta(\tau)|^{-2\kappa}
= Z_1(\cH_\kappa;\,\tau).
\]
The $S$-transformation invariance is automatic from the
Fredholm determinant structure: $\det(1 - \sewop_q)$ depends
only on the eigenvalues, which are determined by the conformal
structure of $E_\tau$, not the choice of fundamental domain.
\end{remark}

\begin{computation}[Genus-1 Fredholm determinant: explicit series;
\ClaimStatusProvedHere]%
\label{comp:thqg-X-genus-1-series}%
\index{Fredholm determinant!genus 1!explicit series}%
The Fredholm determinant at genus~$1$ has the $q$-expansion:
\begin{align}
\det(1 - \sewop_q)^{-1}
&= \prod_{n=1}^{\infty} (1 - q^n)^{-1}
\notag\\
&= 1 + q + 2q^2 + 3q^3 + 5q^4 + 7q^5 + 11q^6 + 15q^7
+ 22q^8 + \cdots
\notag\\
&= \sum_{n=0}^{\infty} p(n)\,q^n,
\label{eq:thqg-X-partition-numbers}
\end{align}
where $p(n)$ is the number of integer partitions of~$n$.
The asymptotic behavior is given by the Hardy--Ramanujan formula:
\begin{equation}\label{eq:thqg-X-hardy-ramanujan}
p(n) \;\sim\;
\frac{1}{4n\sqrt{3}}\,
\exp\!\Bigl(\pi\sqrt{2n/3}\Bigr)
\quad\text{as } n \to \infty.
\end{equation}
The subexponential growth $\log p(n) \sim \pi\sqrt{2n/3}$
is precisely the condition needed for HS-sewing
(Theorem~\ref{thm:general-hs-sewing}).
\end{computation}

\begin{remark}[Trace-class vs.\ Fredholm: comparison of
convergence]%
\label{rem:thqg-X-trace-vs-fredholm}%
The trace $\sum_n q^{n+1} = q/(1-q)$ of $\sewop_q$ converges
for $|q| < 1$, as does every $p$-Schatten norm.
The Fredholm determinant
$\det(1 - \sewop_q)^{-1}
= \prod_n (1 - q^n)^{-1}$
converges absolutely for $|q| < 1$ (the Euler function).
The radius of convergence in~$q$ is $|q| = 1$, where the
product diverges because $\sum_n q^n$ diverges.
On the modular side, $|q| = 1$ corresponds to
$\operatorname{Im}\tau = 0$, the cusp.
\end{remark}

\subsubsection{Genus 2: two handles and the Siegel modular form}

\begin{theorem}[Genus-2 Fredholm determinant;
\ClaimStatusProvedHere]%
\label{thm:thqg-X-genus-2-fredholm}%
\index{Fredholm determinant!genus 2}%
The genus-$2$ partition function of $\cH_\kappa$ on a
surface $\Sigma_2$ with period matrix
$\Omega = \bigl(\begin{smallmatrix}\tau_1 & z \\ z & \tau_2
\end{smallmatrix}\bigr)
\in \Siegel_2$ is
\begin{equation}\label{eq:thqg-X-genus-2-det}
Z_2(\cH_\kappa;\,\Omega)
\;=\;
\bigl(\det\operatorname{Im}\Omega\bigr)^{-\kappa/2}
\cdot
\bigl(\zetareg\Delta_{\Sigma_2}\bigr)^{-\kappa/2}.
\end{equation}
The holomorphic part is the Fredholm determinant of the
genus-$2$ sewing kernel $\sewker_2$ acting on $A^2(D)^{\oplus 2}$:
\begin{equation}\label{eq:thqg-X-genus-2-hol}
Z_2^{\mathrm{hol}}(\cH_1;\,\Omega)
\;=\;
\det\bigl(1 - \sewker_2\bigr)^{-1},
\end{equation}
where $\sewker_2$ is the $2 \times 2$ matrix operator on
$A^2(D)^{\oplus 2}$ with entries determined by the sewing data
of the pair-of-pants decomposition.
\end{theorem}

\begin{proof}
A genus-$2$ surface admits a pair-of-pants decomposition into
$2$ pairs of pants ($2g-2 = 2$) with $3$ sewing circles
($3g-3 = 3$).

\emph{Step 1: pants decomposition.}
Fix a pair-of-pants decomposition: $\Sigma_2 = P_1 \cup_{\partial} P_2$,
where $P_1$ and $P_2$ are each pairs of pants (three-holed spheres).
The three sewing circles have parameters $q_1, q_2, q_3$, where
$q_1$ and $q_2$ correspond to the two handles (B-cycles) and
$q_3$ corresponds to the internal seam.

\emph{Step 2: one-particle operator.}
On the one-particle Bergman space, the sewing produces a
composition of three operators.  Write the one-particle sewing
kernel as a $2 \times 2$ block operator on
$A^2(D)^{\oplus 2}$ (one block per handle):
\begin{equation}\label{eq:thqg-X-genus-2-block}
\sewker_2
=
\begin{pmatrix}
\sewop_{q_1} & R_{12} \\
R_{21} & \sewop_{q_2}
\end{pmatrix},
\end{equation}
where $\sewop_{q_j}$ are the handle sewing operators and
$R_{12}, R_{21}$ are the off-diagonal composition operators
encoding the coupling between the two handles through the internal
seam.  The off-diagonal entries satisfy
$\|R_{12}\|_{\hscl}, \|R_{21}\|_{\hscl} \leq C\,q_3^{1/2}$
by Lemma~\ref{lem:thqg-X-composition-decay}.

\emph{Step 3: Fredholm determinant.}
Since all diagonal and off-diagonal entries are trace class
(the diagonal operators $\sewop_{q_j}$ are trace class by
Lemma~\ref{lem:thqg-X-sewing-schatten}, and the off-diagonal
operators are Hilbert--Schmidt with $\hscl \cdot \hscl \subset \trcl$),
the full operator $\sewker_2$ is trace class on
$A^2(D)^{\oplus 2}$.  The Fredholm determinant is:
\begin{align}
\det(1 - \sewker_2)
&= \det\!\begin{pmatrix}
1 - \sewop_{q_1} & -R_{12} \\
-R_{21} & 1 - \sewop_{q_2}
\end{pmatrix}
\notag\\
&= \det(1 - \sewop_{q_1})\,
\det\bigl(1 - \sewop_{q_2}
- R_{21}(1-\sewop_{q_1})^{-1}R_{12}\bigr).
\label{eq:thqg-X-genus-2-schur}
\end{align}
The Schur complement
$S = \sewop_{q_2} + R_{21}(1-\sewop_{q_1})^{-1}R_{12}$
captures the interaction between the two handles.

\emph{Step 4: degeneration limits.}
In the separating degeneration $q_3 \to 0$
(pinching the internal seam):
$R_{12}, R_{21} \to 0$, and
\[
\det(1 - \sewker_2) \to
\det(1 - \sewop_{q_1}) \cdot \det(1 - \sewop_{q_2})
= \dedeta(\tau_1)^2 \cdot \dedeta(\tau_2)^2,
\]
recovering the product of two genus-$1$ partition functions.
In the non-separating degeneration $q_1 \to 0$
(collapsing one handle):
$\sewop_{q_1} \to 0$ and $\det(1 - \sewker_2)$ reduces to
$\det(1 - \sewop_{q_2} - R_{21}R_{12})$, a genus-$1$
Fredholm determinant with modified sewing operator.

\emph{Step 5: Polyakov formula.}
Applying the Polyakov--Alvarez formula
\eqref{eq:thqg-X-polyakov} at genus~$2$ gives
$\det(1 - \sewker_2)^{-1}
= (\det\operatorname{Im}\Omega)^{-1/2}
\cdot (\zetareg\Delta_{\Sigma_2})^{-1/2}$.
Raising to the power $\kappa$ yields the stated formula.
\end{proof}

\begin{computation}[Genus-2 Fredholm determinant: leading terms;
\ClaimStatusProvedHere]%
\label{comp:thqg-X-genus-2-series}%
\index{Fredholm determinant!genus 2!leading terms}%
Write $q_1 = e^{2\pi i \tau_1}$, $q_2 = e^{2\pi i \tau_2}$,
$r = e^{2\pi i z}$ (the off-diagonal period).
In the product expansion:
\begin{align}
\det(1 - \sewker_2)^{-1}
&= \prod_{n=1}^{\infty}(1-q_1^n)^{-1}
\cdot \prod_{n=1}^{\infty}(1-q_2^n)^{-1}
\cdot \bigl(1 + O(r)\bigr)
\notag\\
&= \frac{1}{\dedeta(\tau_1)\,\dedeta(\tau_2)}
\cdot \Bigl(1 + \sum_{n,m \geq 1}
c_{n,m}\,r^n\,\bar{r}^m + \cdots\Bigr).
\label{eq:thqg-X-genus-2-expansion}
\end{align}
The leading correction $c_{1,1}$ arises from the one-particle
Schur complement and equals
$\sum_{k \geq 0} q_1^{k+1}\,q_2^{k+1}/(1-q_1^{k+1})(1-q_2^{k+1})$.
For $\tau_1 = \tau_2 = i$ (square torus) and $z = i/2$:
$|r| = e^{-\pi}$ and $|\dedeta(i)|^{-2} = \Gamma(1/4)^2/(2\pi^{3/2})$.
\end{computation}

\subsubsection{Genus 3: the three-handle Fredholm determinant}

\begin{theorem}[Genus-3 Fredholm determinant;
\ClaimStatusProvedHere]%
\label{thm:thqg-X-genus-3-fredholm}%
\index{Fredholm determinant!genus 3}%
The genus-$3$ partition function of $\cH_\kappa$ on
$\Sigma_3$ with period matrix $\Omega \in \Siegel_3$ is
\begin{equation}\label{eq:thqg-X-genus-3-det}
Z_3(\cH_\kappa;\,\Omega)
\;=\;
\bigl(\det\operatorname{Im}\Omega\bigr)^{-\kappa/2}
\cdot
\bigl(\zetareg\Delta_{\Sigma_3}\bigr)^{-\kappa/2}.
\end{equation}
The one-particle sewing kernel $\sewker_3$ acts on
$A^2(D)^{\oplus 3}$ as a $3 \times 3$ operator:
\begin{equation}\label{eq:thqg-X-genus-3-block}
\sewker_3
=
\begin{pmatrix}
\sewop_{q_1} & R_{12} & R_{13} \\
R_{21} & \sewop_{q_2} & R_{23} \\
R_{31} & R_{32} & \sewop_{q_3}
\end{pmatrix},
\end{equation}
where the diagonal entries $\sewop_{q_j}$ correspond to the
three handles and the off-diagonal entries $R_{ij}$ encode the
pair-of-pants composition operators for the internal seams.
\end{theorem}

\begin{proof}
A genus-$3$ surface admits a pair-of-pants decomposition into
$4$ pairs of pants ($2g - 2 = 4$) with $6$ sewing circles
($3g - 3 = 6$).  Three of these circles correspond to handles
(B-cycles), and the remaining three correspond to internal
seams.

\emph{Step 1: block structure.}
The $3 \times 3$ block operator $\sewker_3$ on
$A^2(D)^{\oplus 3}$ has diagonal entries $\sewop_{q_j}$
for the three handle sewing operators with
$\|\sewop_{q_j}\|_1 = q_j/(1-q_j)$ (trace class), and
off-diagonal entries $R_{ij}$ with
$\|R_{ij}\|_{\hscl} \leq C\,q_{ij}^{1/2}$
(Hilbert--Schmidt) where $q_{ij}$ is the sewing parameter
of the internal seam connecting handles $i$ and $j$.

\emph{Step 2: trace-class verification.}
The operator $\sewker_3$ is trace class on $A^2(D)^{\oplus 3}$.
The trace norm satisfies
\[
\|\sewker_3\|_1
\leq \sum_{j=1}^{3} \|\sewop_{q_j}\|_1
+ \sum_{i \neq j} \|R_{ij}\|_1
\leq \sum_{j=1}^{3} \frac{q_j}{1-q_j}
+ \sum_{i \neq j}
\|R_{ij}\|_{\hscl}^2
< \infty.
\]

\emph{Step 3: Fredholm determinant.}
The Fredholm determinant is computed by the standard $3 \times 3$
block formula:
\begin{align}
\det(1 - \sewker_3)
&= \det(1 - \sewop_{q_1})
\cdot \det(1 - S_{23|1})
\notag\\
&\quad\text{where } S_{23|1}
= \begin{pmatrix} \sewop_{q_2} + R_{21}(1-\sewop_{q_1})^{-1}R_{12}
& R_{23} + R_{21}(1-\sewop_{q_1})^{-1}R_{13}\\
R_{32} + R_{31}(1-\sewop_{q_1})^{-1}R_{12}
& \sewop_{q_3} + R_{31}(1-\sewop_{q_1})^{-1}R_{13}
\end{pmatrix}
\notag
\end{align}
is the Schur complement.  Iterating the Schur complement
reduction gives:
\begin{equation}\label{eq:thqg-X-genus-3-iterated-schur}
\det(1 - \sewker_3)
= \prod_{j=1}^{3} \det(1 - \sewop_{q_j})
\cdot \det(1 - \Delta_3),
\end{equation}
where $\Delta_3$ is a trace-class operator encoding the
inter-handle coupling.  In the degeneration limit where all
internal seams are pinched ($R_{ij} \to 0$), $\Delta_3 \to 0$
and the determinant factors into a product of three genus-$1$
determinants:
\[
\det(1 - \sewker_3)
\to
\prod_{j=1}^{3}
\dedeta(\tau_j)^2.
\]

\emph{Step 4: Polyakov formula.}
The Polyakov--Alvarez formula at genus~$3$ gives
$\det(1-\sewker_3)^{-1}
= (\det\operatorname{Im}\Omega)^{-1/2}
\cdot (\zetareg\Delta_{\Sigma_3})^{-1/2}$.
\end{proof}

\begin{remark}[The Schottky problem at genus~$3$]%
\label{rem:thqg-X-schottky}%
\index{Schottky problem}%
At genus~$3$, the Siegel upper half-space $\Siegel_3$ has
complex dimension $\binom{3+1}{2} = 6$, and the moduli space
$\cM_3$ has dimension $3g-3 = 6$.  The Torelli map
$\cM_3 \to \Siegel_3/\Sp_6(\Z)$ is therefore generically
surjective at genus~$3$ (this is the classical fact that the
Schottky locus equals the full Siegel space at genus~$\leq 3$).
At genus~$4$, the Siegel space has dimension $10$ while $\cM_4$
has dimension $9$, and the Schottky locus becomes a proper
subvariety.  The Fredholm determinant formula
\eqref{eq:thqg-X-genus-3-det} is therefore a function on all of
$\Siegel_3$ at genus~$3$, but only on the Schottky locus of
$\Siegel_4$ at genus~$4$.
\end{remark}

\subsubsection{General genus: the recursive structure}

\begin{theorem}[General-genus Fredholm structure;
\ClaimStatusProvedHere]%
\label{thm:thqg-X-general-genus-fredholm}%
\index{Fredholm determinant!general genus}%
For every genus $g \geq 1$, the partition function of
$\cH_\kappa$ on $\Sigma_g$ with period matrix
$\Omega \in \Siegel_g$ is
\begin{equation}\label{eq:thqg-X-general-genus-formula}
Z_g(\cH_\kappa;\,\Omega)
\;=\;
\bigl(\det\operatorname{Im}\Omega\bigr)^{-\kappa/2}
\cdot
\bigl(\zetareg\Delta_{\Sigma_g}\bigr)^{-\kappa/2}
\;=\;
\det\bigl(1 - \sewker_g\bigr)^{-\kappa},
\end{equation}
where $\sewker_g$ is a trace-class operator on
$A^2(D)^{\oplus g}$, the genus-$g$ sewing kernel of
Definition~\textup{\ref{def:thqg-X-genus-g-sewing}}.
The following recursive properties hold:
\begin{enumerate}[label=\textup{(\roman*)}]
\item \emph{Separating degeneration.}
  If $\Sigma_g$ degenerates into $\Sigma_{g_1} \cup \Sigma_{g_2}$
  with $g_1 + g_2 = g$, then
  \begin{equation}\label{eq:thqg-X-separating-factorization}
  Z_g(\cH_\kappa) \to
  Z_{g_1}(\cH_\kappa) \cdot Z_{g_2}(\cH_\kappa).
  \end{equation}
\item \emph{Non-separating degeneration.}
  If $\Sigma_g$ degenerates by pinching a non-separating cycle,
  then
  \begin{equation}\label{eq:thqg-X-nonseparating-limit}
  Z_g(\cH_\kappa)
  \;\to\;
  \operatorname{Tr}_{A^2}\bigl(\sewop_q \cdot Z_{g-1}(\cH_\kappa)
  \bigr),
  \end{equation}
  where the trace is over the one-particle space of the
  collapsing handle and $q \to 0$ is the sewing parameter.
\item \emph{Free energy.}
  The genus-$g$ free energy
  $F_g(\cH_\kappa) = \log Z_g(\cH_\kappa)$ satisfies
  \begin{equation}\label{eq:thqg-X-free-energy-genus}
  F_g(\cH_\kappa)
  \;=\;
  -\kappa\,\operatorname{Tr}\log(1 - \sewker_g)
  \;=\;
  \kappa\sum_{k=1}^{\infty}
  \frac{1}{k}\operatorname{Tr}(\sewker_g^k).
  \end{equation}
\end{enumerate}
\end{theorem}

\begin{proof}
The general formula follows from the Heisenberg sewing theorem
(Theorem~\ref{thm:thqg-X-heisenberg-sewing-full}\,(IV)) and
the Polyakov--Alvarez formula at genus~$g$.

Part~(i): in the separating degeneration, the genus-$g$
sewing kernel $\sewker_g$ becomes block-diagonal:
$\sewker_g \to \sewker_{g_1} \oplus \sewker_{g_2}$
as the separating seam parameter $q_s \to 0$.
Hence $\det(1 - \sewker_g) \to
\det(1 - \sewker_{g_1}) \cdot \det(1 - \sewker_{g_2})$.

Part~(ii): in the non-separating degeneration, one diagonal
block $\sewop_{q_j} \to 0$ as $q_j \to 0$, and the Schur
complement formula reduces the rank of the sewing kernel by one.

Part~(iii): $\log\det(1-\sewker_g)^{-\kappa}
= -\kappa\,\operatorname{Tr}\log(1-\sewker_g)
= \kappa\sum_{k} k^{-1}\operatorname{Tr}(\sewker_g^k)$.
\end{proof}

\begin{proposition}[Free energy: $\hat{A}$-genus verification;
\ClaimStatusProvedHere]%
\label{prop:thqg-X-free-energy-ahat}%
\index{free energy!$\hat{A}$-genus}%
For $\cH_\kappa$ at level $\kappa$, the Faber--Pandharipande
free energy formula
$F_g = \kappa \cdot \lambda_g^{\mathrm{FP}}$ is verified
through $g = 3$:
\begin{center}
\begin{tabular}{c|c|c|c}
$g$ & $\lambda_g^{\mathrm{FP}}$ & $F_g/\kappa$ & Source \\
\hline
$1$ & $\frac{|B_2|}{2 \cdot 2!} = \frac{1}{24}$ & $\frac{1}{24}$
& $\dedeta(\tau)^{-1}$: constant term \\[4pt]
$2$ & $\frac{2^3-1}{2^3}\frac{|B_4|}{4!}
= \frac{7}{8}\cdot\frac{1}{720}
= \frac{7}{5760}$ & $\frac{7}{5760}$
& Polyakov at $g = 2$ \\[4pt]
$3$ & $\frac{2^5-1}{2^5}\frac{|B_6|}{6!}
= \frac{31}{32}\cdot\frac{1}{30240}
= \frac{31}{967680}$ & $\frac{31}{967680}$
& Polyakov at $g = 3$
\end{tabular}
\end{center}
These match the $\hat{A}$-genus coefficients
$\hat{A}_g = \frac{(2^{2g-1}-1)|B_{2g}|}{2^{2g-1}(2g)!}$
(Theorem~\textup{\ref{thm:family-index}}).
\end{proposition}

\begin{proof}
The genus-$g$ free energy of a single free boson is
$F_g = -\frac{1}{2}\log\zetareg\Delta_{\Sigma_g}$
integrated over $\overline{\cM}_g$.  By Faltings'
arithmetic Riemann--Roch formula and the Mumford relations,
$\int_{\overline{\cM}_g} \lambda_g
= \frac{(2^{2g-1}-1)|B_{2g}|}{2^{2g-1}(2g)!}$
(Faber's formula).  The claim follows by direct substitution.
\end{proof}

% =====================================================================
%
%  SECTION 4: SCALAR SATURATION AND CLASS-G ALGEBRAS
%
% =====================================================================

\subsection{Extension to class-G algebras via scalar saturation}
\label{subsec:thqg-X-class-G}

The Heisenberg sewing theorem gives a Fredholm determinant
formula for the simplest possible chiral algebra, the free boson.
The shadow depth classification
(Definition~\ref{def:shadow-depth-classification}
in Appendix~\ref{app:nonlinear-modular-shadows})
organizes chiral algebras into four classes based on the
termination behavior of the shadow obstruction tower:
\begin{itemize}
\item Class~\textbf{G} (Gaussian): $r_{\max} = 2$.
  The shadow obstruction tower terminates at weight~$2$:
  $\Theta_\cA^{\min} = \eta\otimes\Gamma_\cA$.
  On the proved uniform-weight scalar lane this specializes to
  $\Theta_\cA^{\min} = \kappa(\cA)\cdot\eta \otimes \Lambda$.
  Examples: $\cH_\kappa$, lattice algebras $V_\Lambda$.
\item Class~\textbf{L} (Lie): $r_{\max} = 3$.
  Examples: $V_k(\fg)$ at generic level.
\item Class~\textbf{C} (contact): $r_{\max} = 4$.
  Examples: $\beta\gamma$ system.
\item Class~\textbf{M} (mixed): $r_{\max} = \infty$.
  Examples: $\mathrm{Vir}_c$, $\cW_N$ for $N \geq 3$.
\end{itemize}
For class-G algebras, the scalar rigidity theorem
(Theorem~\ref{thm:algebraic-family-rigidity}) ensures only that the
minimal MC class is purely scalar in arity.  This does not by itself
identify the full nonperturbative partition function with a Fredholm
determinant.  The Fredholm formula below is proved for Gaussian
algebras on the uniform-weight scalar lane.

\begin{theorem}[Fredholm scalar partition function on the Gaussian
scalar lane; \ClaimStatusProvedHere]%
\label{thm:thqg-X-class-G-fredholm}%
\index{partition function!class G}%
\index{scalar saturation!Fredholm determinant}%
Let $\cA$ be a class-G Gaussian chiral algebra on the proved
uniform-weight scalar lane
\textup{(}shadow depth $r_{\max} = 2$\textup{)} with
modular characteristic $\kappa = \kappa(\cA)$ and central
charge $c(\cA)$.  Assume HS-sewing
\textup{(}Definition~\textup{\ref{def:hs-sewing}}\textup{)}.
Then the scalar genus-$g$ partition function is a Fredholm
determinant:
\begin{equation}\label{eq:thqg-X-class-G-formula}
Z_g^{\mathrm{scal}}(\cA;\,\Omega)
\;=\;
Z_g^{\mathrm{G}}(\kappa;\,\Omega)
\;:=\;
\bigl(\det\operatorname{Im}\Omega\bigr)^{-\kappa/2}
\cdot
\bigl(\zetareg\Delta_{\Sigma_g}\bigr)^{-\kappa/2},
\end{equation}
up to a $\kappa$-dependent prefactor encoding the vacuum energy
$q^{-c/24}$.  This formula depends on $\cA$ only through
$\kappa(\cA)$ on that scalar lane.  For general multi-weight
class-G algebras, the identification of the scalar coefficient with
$\kappa(\cA)\Lambda$ remains open, so no full Fredholm theorem is
claimed here.
\end{theorem}

\begin{proof}
On the uniform-weight Gaussian lane,
Theorem~\ref{thm:theta-direct-derivation} identifies the minimal
scalar package as
$\Theta_\cA^{\min} = \kappa(\cA)\eta\otimes\Lambda$.
The scalar genus-$g$ partition function is therefore the
$\kappa(\cA)$-multiple of the Heisenberg scalar partition function.
The HS-sewing condition ensures convergence of the Fredholm
determinant (by Theorem~\ref{thm:general-hs-sewing}, which
requires polynomial OPE growth and subexponential sector growth;
class-G algebras satisfy both).

The Fredholm determinant formula then follows from the
Heisenberg computation
(Theorem~\ref{thm:thqg-X-heisenberg-sewing-full}\,(IV))
with $\kappa$ replaced by $\kappa(\cA)$.
\end{proof}

\begin{example}[Lattice vertex algebras]%
\label{ex:thqg-X-lattice-fredholm}%
\index{lattice algebra!Fredholm determinant}%
Let $V_\Lambda$ be a lattice vertex algebra associated to an
even integral lattice $\Lambda$ of rank $r$.  Then
$\kappa(V_\Lambda) = r$ (the rank of the lattice), and the
genus-$g$ partition function is
\begin{equation}\label{eq:thqg-X-lattice-formula}
Z_g(V_\Lambda;\,\Omega)
\;=\;
\bigl(\det\operatorname{Im}\Omega\bigr)^{-r/2}
\cdot
\bigl(\zetareg\Delta_{\Sigma_g}\bigr)^{-r/2}
\cdot
\Theta_\Lambda^{(g)}(\Omega),
\end{equation}
where $\Theta_\Lambda^{(g)}(\Omega)
= \sum_{\lambda \in \Lambda^g}
\exp(\pi i \lambda^T \Omega \lambda)$ is the genus-$g$
lattice theta function.  The Fredholm part (first two factors)
carries the oscillator modes and is universal in $r$;
the theta function carries the zero-mode sum over the
lattice and is lattice-specific.

The HS-sewing condition is verified: $V_\Lambda$ has sector
growth $\dim (V_\Lambda)_n \leq |\Lambda^*/\Lambda| \cdot p(n)^r$
(polynomial times partition function growth), and OPE
coefficients are bounded by $\kappa = r$ times the lattice
norm (Corollary~\ref{cor:hs-sewing-standard-landscape}).
\end{example}

\begin{remark}[Curvature-braiding orthogonality for lattice algebras]%
\label{rem:thqg-X-lattice-curvature-braiding}%
\index{lattice algebra!curvature-braiding orthogonality}%
By Theorem~\ref{thm:lattice:curvature-braiding-orthogonal},
$\kappa(V_\Lambda^{N,q}) = \mathrm{rank}(\Lambda)$
independent of the cocycle~$q$.  This means the Fredholm
determinant part of the partition function is insensitive to
the choice of cocycle; only the theta function
$\Theta_\Lambda^{(g)}$ depends on the lattice embedding.
This orthogonality between the oscillator (Fredholm) sector
and the zero-mode (theta) sector is a direct consequence of
the class-G property.
\end{remark}

\begin{proposition}[Rank additivity;
\ClaimStatusProvedHere]%
\label{prop:thqg-X-rank-additivity}%
\index{Fredholm determinant!rank additivity}%
If $\cA = \cA_1 \otimes \cA_2$ is a tensor product of
class-G algebras, then
\begin{equation}\label{eq:thqg-X-rank-additivity}
Z_g(\cA;\,\Omega)
= Z_g(\cA_1;\,\Omega)
\cdot Z_g(\cA_2;\,\Omega),
\end{equation}
reflecting the additivity of the modular characteristic
$\kappa(\cA) = \kappa(\cA_1) + \kappa(\cA_2)$
and the multiplicativity of the Fredholm determinant
$\det(1 - T_1 \oplus T_2)
= \det(1 - T_1) \cdot \det(1 - T_2)$.
\end{proposition}

\begin{proof}
By Proposition~\ref{prop:independent-sum-factorization},
the modular characteristic is additive:
$\kappa(\cA_1 \otimes \cA_2)
= \kappa(\cA_1) + \kappa(\cA_2)$.
The sewing operator for a tensor product is the direct sum:
$\sewker_g(\cA_1 \otimes \cA_2)
= \sewker_g(\cA_1) \oplus \sewker_g(\cA_2)$
on $A^2(D)^{\oplus g_1} \oplus A^2(D)^{\oplus g_2}$.
The Fredholm determinant of a direct sum is the product:
$\det(1 - T_1 \oplus T_2) = \det(1-T_1)\det(1-T_2)$.
\end{proof}

% =====================================================================
%
%  SECTION 5: NON-FREDHOLM CORRECTIONS FOR CLASSES L, C, M
%
% =====================================================================

\subsection{Non-Fredholm corrections for classes L, C, M}
\label{subsec:thqg-X-non-fredholm}

For algebras of shadow depth $r_{\max} > 2$, the partition
function is no longer a pure Fredholm determinant.  The higher
shadow data ($\fC$, $\mathfrak{Q}$, and all higher
$o_r$ for $r \geq 3$) produce perturbative corrections
to the Gaussian (Fredholm) base, organized as Feynman
integrals weighted by the graph-sum formula of
Definition~\ref{def:modular-bar-hamiltonian}.

\begin{definition}[Non-Fredholm correction]%
\label{def:thqg-X-non-fredholm-correction}%
\index{non-Fredholm correction|textbf}%
\index{partition function!non-Fredholm correction}%
For a chiral algebra $\cA$ with shadow depth $r_{\max}$, define
the \emph{genus-$g$ non-Fredholm correction} as
\begin{equation}\label{eq:thqg-X-non-fredholm-correction}
\delta Z_g(\cA;\,\Omega)
\;:=\;
Z_g(\cA;\,\Omega)
\;-\;
Z_g^{\mathrm{G}}(\kappa(\cA);\,\Omega),
\end{equation}
where $Z_g^{\mathrm{G}}$ is the Gaussian (Fredholm) base of
\eqref{eq:thqg-X-class-G-formula}.  The correction
$\delta Z_g$ encodes the contribution of all shadow data
beyond the modular characteristic $\kappa$.
\end{definition}

\begin{theorem}[Feynman integral expansion;
\ClaimStatusProvedHere]%
\label{thm:thqg-X-feynman-expansion}%
\index{Feynman integral!non-Fredholm correction}%
For a modular Koszul chiral algebra $\cA$ satisfying HS-sewing,
the genus-$g$ non-Fredholm correction admits a Feynman integral
expansion:
\begin{equation}\label{eq:thqg-X-feynman-expansion}
\delta Z_g(\cA;\,\Omega)
\;=\;
\sum_{r=3}^{r_{\max}}
\sum_{\Gamma \in \mathcal{G}^{\mathrm{conn}}_{g,r}}
\frac{1}{|\Aut(\Gamma)|}
\int_{\Sigma_g^{v(\Gamma)}}
\omega_\Gamma(\Omega),
\end{equation}
where:
\begin{itemize}
\item $\mathcal{G}^{\mathrm{conn}}_{g,r}$ is the set of
  connected stable graphs of genus~$g$ with vertices of
  arity at most~$r$;
\item $v(\Gamma)$ is the number of vertices of~$\Gamma$;
\item $\omega_\Gamma(\Omega)$ is the integrand determined by
  the graph~$\Gamma$, with vertex labels from the transferred
  cyclic minimal model $\ell_k^{\mathrm{tr}}$ of $\cA$ and
  edge contractions by the complementarity propagator
  $P_\cA = H_\cA^{-1}$ (the inverse of the Hodge Laplacian).
\end{itemize}
The Feynman expansion converges absolutely for $|q_\alpha| < 1$
(all sewing parameters inside the unit disk).
\end{theorem}

\begin{proof}
The genus-$g$ partition function is computed by the shadow
Postnikov tower: $Z_g = \langle e^{\Theta_\cA^{(g)}} \rangle_g$.
The Gaussian part $Z_g^{\mathrm{G}}$ is the contribution of
$\Theta_\cA^{\leq 2}$, which consists of the modular
characteristic $\kappa$ contracted with the Hodge cycle.
The non-Fredholm correction is the contribution of
$\Theta_\cA - \Theta_\cA^{\leq 2}$, which involves the
higher-arity terms $\ell_k^{\mathrm{tr}}$ for $k \geq 3$.

Expanding $e^{\Theta}$ in the graph-sum formula
(Definition~\ref{def:modular-bar-hamiltonian}), each
graph $\Gamma$ with a vertex of arity $r > 2$ contributes to
$\delta Z_g$.  The weight of a graph is
$w(\Gamma) = \sum_{v \in V(\Gamma)} \mathrm{arity}(v)$,
and graphs with $w > 2|V(\Gamma)|$ contribute to the
non-Fredholm correction.

Convergence: the HS-sewing condition
(Theorem~\ref{thm:general-hs-sewing}) ensures that each
Feynman integral converges absolutely.  The propagator
$P_\cA = H_\cA^{-1}$ on $\Sigma_g$ has the same exponential
decay as the Bergman composition operator
(Lemma~\ref{lem:thqg-X-composition-decay}), and the vertex
labels $\ell_k^{\mathrm{tr}}$ are bounded by the polynomial
OPE growth.  The arity-cutoff lemma (Lemma~\ref{lem:arity-cutoff})
bounds the number of contributing graphs at each order.
\end{proof}

\subsubsection{Class L: cubic corrections from Lie algebras}

\begin{proposition}[Class-L Feynman integrals;
\ClaimStatusProvedHere]%
\label{prop:thqg-X-class-L-feynman}%
\index{class L!Feynman integrals}%
For a class-L algebra $\cA$ \textup{(}$r_{\max} = 3$,
e.g.\ $V_k(\fg)$ at generic level\textup{)}, the
non-Fredholm correction at genus~$g$ involves only
trivalent graphs:
\begin{equation}\label{eq:thqg-X-class-L-correction}
\delta Z_g^{\mathrm{L}}(\cA;\,\Omega)
\;=\;
\sum_{\Gamma \in \mathcal{G}^{3\text{-val}}_{g}}
\frac{1}{|\Aut(\Gamma)|}
\int_{\Sigma_g^{v(\Gamma)}}
\omega_\Gamma^{(3)}(\Omega),
\end{equation}
where $\mathcal{G}^{3\text{-val}}_g$ denotes the set of
connected $3$-valent stable graphs of genus~$g$, and
$\omega_\Gamma^{(3)}$ is the integrand with cubic vertex
labels $\ell_3^{\mathrm{tr}}$ (the transferred $L_\infty$
ternary bracket).

A $3$-valent graph of genus~$g$ with $v$ vertices has
$e = 3v/2$ edges and loop number $g = e - v + 1 = v/2 + 1$,
so $v = 2(g-1)$ and $e = 3(g-1)$.  In particular:
\begin{center}
\begin{tabular}{c|c|c|c}
$g$ & $v$ & $e$ & \# connected $3$-valent graphs \\
\hline
$1$ & $0$ & $0$ & $0$ (no correction) \\
$2$ & $2$ & $3$ & $2$ (theta, sunset) \\
$3$ & $4$ & $6$ & \text{finitely many}
\end{tabular}
\end{center}
The first non-trivial correction appears at genus~$2$
with the theta graph and the sunset graph.
\end{proposition}

\begin{proof}
For class-L algebras, $\ell_k^{\mathrm{tr}} = 0$ for $k \geq 4$
(the shadow obstruction tower terminates at arity~$3$).  Hence only graphs
with vertices of arity exactly~$3$ contribute to $\delta Z_g$.
These are the $3$-regular (trivalent) graphs.

The graph count follows from the Euler relation: a connected
graph with $v$ vertices, $e$ edges, and first Betti number
$g$ satisfies $g = e - v + 1$.  For $3$-regular graphs,
$3v = 2e$ (every edge has two endpoints; every vertex has
three), so $e = 3v/2$ and $g = v/2 + 1$.
\end{proof}

\begin{example}[Affine $\AffKM{g}_k$ at genus~$2$]%
\label{ex:thqg-X-affine-genus-2}%
\index{affine Kac--Moody!genus-2 correction}%
For $V_k(\fg)$ at generic level $k \neq -h^\vee$, the
genus-$2$ non-Fredholm correction involves two Feynman graphs.
Let $f_{abc}$ denote the structure constants of $\fg$
and $P(z,w;\Omega)$ the genus-$2$ propagator.  Then:
\begin{enumerate}[label=\textup{(\alph*)}]
\item \emph{Theta graph}
  $(\Gamma_\theta$: two vertices, three edges$)$:
  \begin{equation}\label{eq:thqg-X-theta-graph}
  I_{\theta}
  = \frac{1}{12}
  f_{abc}\,f_{a'b'c'}\,
  \langle a,a'\rangle\langle b,b'\rangle\langle c,c'\rangle
  \int_{\Sigma_2 \times \Sigma_2}
  P(z,w;\Omega)^3\,dz\,dw.
  \end{equation}
\item \emph{Sunset graph}
  $(\Gamma_{\mathrm{sun}}$: two vertices, three edges,
  one edge connects vertex to itself$)$:
  \begin{equation}\label{eq:thqg-X-sunset-graph}
  I_{\mathrm{sun}}
  = \frac{1}{4}
  f_{abc}\,f_{a'b'c'}\,
  \langle a,a'\rangle
  \int_{\Sigma_2}
  \langle b, c\rangle_z
  \,P(z,w;\Omega)\,
  \langle b', c'\rangle_w\,
  dz\,dw.
  \end{equation}
\end{enumerate}
The total correction is
$\delta Z_2^{\mathrm{L}}(V_k(\fg);\,\Omega)
= I_\theta + I_{\mathrm{sun}}$.
At integrable level $k \in \Z_{\geq 0}$, these Feynman
integrals reduce to combinatorial sums via the Verlinde
formula (Example~\ref{eq:verlinde-general}).
\end{example}

\subsubsection{Class C: quartic contact corrections}

\begin{proposition}[Class-C quartic contact;
\ClaimStatusProvedHere]%
\label{prop:thqg-X-class-C-quartic}%
\index{class C!quartic contact}%
For a class-C algebra $\cA$ with $r_{\max} = 4$
\textup{(}e.g.\ the $\beta\gamma$ system\textup{)},
the non-Fredholm correction at genus~$g$ involves
graphs with vertices of arity $3$ and $4$.
The quartic vertex $\ell_4^{\mathrm{tr}}$ is governed by the
quartic resonance class
$\mathfrak{Q} \in H^2(\cA^{\mathrm{sh}}_{4,0})$
of Appendix~\textup{\ref{app:nonlinear-modular-shadows}}.

The first quartic contribution appears at genus~$1$:
\begin{equation}\label{eq:thqg-X-quartic-genus-1}
\delta Z_1^{(4)}(\cA;\,\tau)
\;=\;
\mathfrak{Q}(\cA)\,
\int_{E_\tau}
P(z;\tau)^2\,dz,
\end{equation}
where $P(z;\tau) = G_\tau(z)$ is the genus-$1$ propagator.
For the $\beta\gamma$ system,
$\mathfrak{Q}_{\beta\gamma} = 0$
\textup{(}Corollary~\textup{\ref{cor:nms-betagamma-mu-vanishing}}\textup{)},
so the quartic contact invariant vanishes and the first
genuine quartic correction appears at genus~$2$.
\end{proposition}

\begin{proof}
The quartic vertex contributes Feynman graphs with at least one
$4$-valent vertex.  At genus~$1$, the simplest such graph is a
single $4$-valent vertex with two self-loops, contributing
$\delta Z_1^{(4)}$.  The quartic resonance class $\mathfrak{Q}$
is the weight-$4$ component of the shadow algebra $\cA^{\mathrm{sh}}$
(Corollary~\ref{cor:shadow-extraction}).

For the $\beta\gamma$ system: the quartic contact invariant
$\mu_{\beta\gamma} = \langle \eta, m_3(\eta,\eta,\eta)\rangle = 0$
by Corollary~\ref{cor:nms-betagamma-mu-vanishing} (rank-one abelian
rigidity on the weight-changing line).
\end{proof}

\subsubsection{Class M: the infinite shadow obstruction tower}

\begin{remark}[Class-M non-Fredholm corrections]%
\label{rem:thqg-X-class-M-corrections}%
\index{class M!non-Fredholm corrections}%
For class-M algebras ($r_{\max} = \infty$, e.g.\ $\mathrm{Vir}_c$
and $\cW_N$ for $N \geq 3$), the Feynman expansion involves
vertices of arbitrarily high arity.  The genus-$g$ non-Fredholm
correction is an infinite series:
\[
\delta Z_g^{\mathrm{M}}(\cA;\,\Omega)
= \sum_{r=3}^{\infty}
\sum_{\Gamma \in \mathcal{G}^{\mathrm{conn}}_{g,r}}
\frac{1}{|\Aut(\Gamma)|}
\int_{\Sigma_g^{v(\Gamma)}}
\omega_\Gamma(\Omega).
\]
Convergence of this series is guaranteed by the HS-sewing
condition (Theorem~\ref{thm:general-hs-sewing}), which provides
uniform bounds on the vertex labels $\ell_k^{\mathrm{tr}}$, and
the arity-cutoff lemma (Lemma~\ref{lem:arity-cutoff}), which
ensures that the contribution of high-arity graphs is
exponentially suppressed.

For the Virasoro algebra at genus~$1$, the non-Fredholm
correction involves the quartic contact invariant
$\mathfrak{Q}_{\mathrm{Vir}} = 10/[c(5c+22)]$
and all higher obstruction classes $o_r$ for $r \geq 5$
(the quintic obstruction $o_5$ is forced to be nonzero
by Theorem~\ref{thm:nms-virasoro-quintic-forced}).  The full partition
function $Z_1(\mathrm{Vir}_c;\,\tau)$ at genus~$1$ is:
\[
Z_1(\mathrm{Vir}_c;\,\tau)
= \dedeta(\tau)^{-(c-1)}
= q^{-(c-1)/24}\prod_{n=2}^{\infty}(1-q^n)^{-1},
\]
the character of the Virasoro vacuum module, which is
\emph{not} a pure Fredholm determinant on a single Bergman space
(it is missing the $n=1$ factor from the product).
\end{remark}

\begin{proposition}[Virasoro: Fredholm base + corrections;
\ClaimStatusProvedHere]%
\label{prop:thqg-X-virasoro-decomposition}%
\index{Virasoro!Fredholm decomposition}%
The genus-$1$ Virasoro partition function admits the decomposition
\begin{equation}\label{eq:thqg-X-virasoro-decomposition}
Z_1(\mathrm{Vir}_c;\,\tau)
\;=\;
\underbrace{|\dedeta(\tau)|^{-2c}}_{\text{Gaussian base}
(\kappa = c/2)}
\;\cdot\;
\underbrace{|\dedeta(\tau)|^{2}
\cdot (1-q)}_{\text{null-state correction}},
\end{equation}
where the Gaussian base is the Fredholm determinant at
$\kappa = c/2$ and the correction factor removes the
contribution of the $L_{-1}$ null state.
\end{proposition}

\begin{proof}
The Virasoro character at genus~$1$ is
$\chi_c(\tau) = q^{(1-c)/24}/(1-q)\prod_{n=2}^{\infty}(1-q^n)^{-1}
= q^{-c/24}\,\dedeta(\tau)^{-1}\cdot(1-q)^{-1}\cdot\dedeta(\tau)/q^{-1/24}$\,.
More directly:
\begin{align}
\chi_c(\tau)
&= q^{(1-c)/24} \prod_{n=2}^{\infty}(1-q^n)^{-1}
\notag\\
&= q^{(1-c)/24}
\cdot \frac{\prod_{n=1}^{\infty}(1-q^n)^{-1}}{(1-q)^{-1}}
\notag\\
&= q^{(1-c)/24}
\cdot (1-q) \cdot \prod_{n=1}^{\infty}(1-q^n)^{-1}
\notag\\
&= (1-q)\cdot q^{-c/24}\cdot q^{1/24}
\cdot \prod_{n=1}^{\infty}(1-q^n)^{-1}
\notag\\
&= (1-q)\cdot q^{-c/24}\cdot \dedeta(\tau)^{-1}.
\label{eq:thqg-X-virasoro-factorization}
\end{align}
The factor $(1-q)$ removes the $L_{-1}$ descendant (which is
null for the vacuum module), and $\dedeta(\tau)^{-1}$ is the
oscillator contribution $\prod_{n=1}^{\infty}(1-q^n)^{-1}$
up to the $q^{1/24}$ prefactor.  The Gaussian base
$|\dedeta|^{-2c}$ encodes the oscillator contribution
with modular characteristic $\kappa = c/2$
(Theorem~\ref{thm:modular-characteristic}).
\end{proof}

\begin{remark}[Resurgent structure for class-M algebras]%
\label{rem:thqg-X-resurgent}%
\index{resurgence!class M}%
For class-M algebras, the Feynman expansion
$\delta Z_g^{\mathrm{M}}$ is expected to be resurgent:
the asymptotic (divergent) perturbative series in the shadow
tower is Borel summable, with non-perturbative corrections
controlled by the resonance rank $\rho(\cA)$
(Definition~\ref{def:resonance-rank}).  The Stokes phenomenon
as $\tau$ crosses walls of marginal stability corresponds to
the wall-crossing of the MC moduli of
$\Theta_\cA \in \MC(\gAmod)$.  This resurgent structure is
conjectural beyond the Virasoro case, where the Zamolodchikov
recursion provides an independent control.
\end{remark}

% =====================================================================
%
%  SECTION 6: ANALYTIC ASPECTS
%
% =====================================================================

\subsection{Analytic aspects: the analytic bar coalgebra,
coderived categories, and open problems}
\label{subsec:thqg-X-analytic-aspects}

\subsubsection{The analytic bar coalgebra}

\begin{definition}[Analytic bar coalgebra; review]%
\label{def:thqg-X-analytic-bar-review}%
\index{bar coalgebra!analytic!review}%
Recall from Definition~\ref{def:analytic-bar-coalgebra} that the
analytic bar coalgebra $\anbar(\cA)$ is defined by closing
the algebraic bar construction inside the sewing envelope:
\[
\anbar_n(\cA)
\;=\;
\overline{
\algcore^{\otimes n}
\otimes \Omega^*_{\log}(C_n(X))
}^{\,\graphnorm(d_{\mathrm{bar}})},
\]
where the graph norm of $d_{\mathrm{bar}}$ is the norm
$\|x\|_{\mathrm{graph}}^2 := \|x\|^2 + \|d_{\mathrm{bar}} x\|^2$.
This is the analytic closure of the algebraic bar construction
(Definition~\ref{def:bar-differential-complete}).
\end{definition}

\begin{proposition}[Analytic bar differential is bounded;
\ClaimStatusProvedHere]%
\label{prop:thqg-X-analytic-bar-bounded}%
\index{bar coalgebra!analytic!boundedness}%
Under HS-sewing \textup{(}Definition~\textup{\ref{def:hs-sewing}}\textup{)},
the bar differential $d_{\mathrm{bar}} \colon \anbar_n(\cA) \to
\anbar_{n-1}(\cA)$ is a bounded operator for the graph-norm
topology, with operator norm
\begin{equation}\label{eq:thqg-X-bar-bounded}
\|d_{\mathrm{bar}}\|_{\mathrm{op}}
\;\leq\;
\binom{n}{2}\,K\,\sup_q \|m_q\|_{\hscl},
\end{equation}
where $K$ is the polynomial OPE growth constant and $m_q$ is the
$q$-weighted multiplication.
\end{proposition}

\begin{proof}
The bar differential $d_{\mathrm{bar}} = \sum_{i<j}
\operatorname{Res}_{z_i=z_j} \circ Y_{ij}$ has
$\binom{n}{2}$ summands, each of which is a composition of
the OPE operator $Y_{ij}$ (bounded by $K$ on the sewing envelope)
and the residue operator (norm $1$ on the graph-norm closure).
The HS-sewing condition ensures $\|m_q\|_{\hscl} < \infty$,
which controls the operator norm of $Y_{ij}$ on the weighted
completion $\sewenv_q$.
\end{proof}

\begin{proposition}[Analytic coproduct;
\ClaimStatusProvedHere]%
\label{prop:thqg-X-analytic-coproduct}%
\index{bar coalgebra!analytic!coproduct}%
The analytic bar coalgebra $\anbar(\cA)$ carries a continuous
coproduct
\begin{equation}\label{eq:thqg-X-analytic-coproduct}
\bar{\Delta}^{\mathrm{an}} \colon
\anbar(\cA)
\;\to\;
\anbar(\cA)
\;\widehat{\otimes}\;
\anbar(\cA),
\end{equation}
where $\widehat{\otimes}$ is the completed tensor product
in the graph-norm topology.  The coproduct extends the algebraic
coproduct $\bar{\Delta}$ of the algebraic bar construction
by continuity.
\end{proposition}

\begin{proof}
The algebraic coproduct $\bar{\Delta}$ on $\barB_n(\cA)$
sends $\barB_n \to \bigoplus_{p+q=n} \barB_p \otimes \barB_q$
by deshuffle.  Each deshuffle component is a restriction map
(restricting meromorphic forms from $C_n$ to
$C_p \times C_q$), which is continuous for the graph-norm
topology.  The completed coproduct is the closure of the
algebraic coproduct.
\end{proof}

\subsubsection{Coderived categories and genus $g \geq 1$}

\begin{remark}[Curvature forces coderived passage; review]%
\label{rem:thqg-X-coderived-review}%
\index{coderived category!Fredholm context}%
At genus~$g \geq 1$, the fiberwise bar differential satisfies
$\dfib^{\,2} = \kappa(\cA) \cdot \omega_g \neq 0$
(the curvature of the Hodge connection;
Theorem~\ref{thm:genus-graded-bar}).  This means that the
bar complex at genus~$g$ is \emph{curved}: $d^2 \neq 0$,
and ordinary cohomology is not defined.  The correct
receptacle is the \emph{coderived category}
$\codercat(\anbar_g(\cA)\text{-comod})$
(Definition~\ref{def:analytic-shadow} and
Remark~\ref{rem:curvature-coderived-analytic}), not the
ordinary derived category.

For the Fredholm partition functions of this section, the
coderived passage manifests as follows: the genus-$g$
partition function $Z_g$ is the \emph{supertrace} of the
sewing operator on the coderived category of the analytic bar
coalgebra, not an ordinary cohomological invariant.  The
Fredholm determinant formula
$Z_g = \det(1 - \sewker_g)^{-\kappa}$ computes this
supertrace directly via the one-particle reduction, bypassing
the coderived formalism for the Heisenberg and class-G cases.
For classes L, C, M, the coderived category is essential: the
non-Fredholm corrections $\delta Z_g$ are coderived
invariants that vanish in the ordinary derived category.
\end{remark}

\begin{definition}[Coderived shadow partition function]%
\label{def:thqg-X-coderived-shadow-pf}%
\index{partition function!coderived shadow}%
For a chiral algebra $\cA$ satisfying HS-sewing, define the
\emph{coderived shadow partition function} at genus~$g$ as
\begin{equation}\label{eq:thqg-X-coderived-shadow-pf}
Z_g^{\mathrm{co}}(\cA;\,\tau)
\;=\;
\operatorname{Str}_{Q_g^{\mathrm{an}}(\cA)}
\!\bigl(\sigma_g\,q^{L_0 - c/24}\bigr),
\end{equation}
where $Q_g^{\mathrm{an}}(\cA)
= \codercat(\anbar_g(\cA)\text{-comod})$ is the
genus-$g$ analytic shadow
\textup{(}Definition~\textup{\ref{def:analytic-shadow}}\textup{)},
$\sigma_g$ is the complementarity involution, and
$\operatorname{Str}$ is the supertrace.
\end{definition}

\begin{proposition}[Coderived = Fredholm for class~G;
\ClaimStatusProvedHere]%
\label{prop:thqg-X-coderived-fredholm-G}%
\index{coderived category!Fredholm for class G}%
For class-G algebras, the coderived shadow partition function
coincides with the Fredholm partition function:
\begin{equation}\label{eq:thqg-X-coderived-fredholm-agreement}
Z_g^{\mathrm{co}}(\cA;\,\tau)
\;=\;
Z_g(\cA;\,\Omega)
\;=\;
\det(1 - \sewker_g)^{-\kappa(\cA)}.
\end{equation}
\end{proposition}

\begin{proof}
For class-G algebras, the curvature is scalar
($\dfib^2 = \kappa \cdot \omega_g$), and the coderived category
reduces to the ordinary derived category twisted by the Hodge
line bundle $\hodgebdle^{\kappa/2}$.  The supertrace on
the twisted derived category equals the Fredholm determinant
of the sewing kernel, by the Heisenberg sewing theorem
(Theorem~\ref{thm:thqg-X-heisenberg-sewing-full}).
\end{proof}

\subsubsection{The analytic Koszul pair}

\begin{remark}[Analytic Koszul duality]%
\label{rem:thqg-X-analytic-koszul}%
\index{Koszul pair!analytic}%
An analytic Koszul pair $(\cA, \cA^!)$ consists of an
algebraic Koszul pair equipped with compatible sewing
envelopes and analytic bar coalgebras
(Definition~\ref{def:analytic-koszul-pair}).  For the
Heisenberg algebra, the analytic Koszul pair is
$(\cH_\kappa, \cH_\kappa^!)$,
where $\cH_\kappa^! = \operatorname{Sym}^{\mathrm{ch}}(V^*)$ has
modular characteristic $\kappa(\cH_\kappa^!) = -\kappa$.
The analytic bar coalgebra of $\cH_\kappa^!$
is $\anbar(\cH_\kappa^!) \cong
\overline{\operatorname{Sym} A^2(D)}^{\mathsf{co}}$, the
cofree coalgebra completion of the Bergman Fock space.

The genus-$g$ complementarity decomposition
$Z_g(\cH_\kappa) + Z_g(\cH_\kappa^!)
= (\det\operatorname{Im}\Omega)^{-\kappa/2}
\cdot (\zetareg\Delta)^{-\kappa/2}
+ (\det\operatorname{Im}\Omega)^{\kappa/2}
\cdot (\zetareg\Delta)^{\kappa/2}$
is a direct Fredholm manifestation of Theorem~C
(complementarity): what $\cH_\kappa$ sees as the positive
curvature $+\kappa \cdot \omega_g$, its dual $\cH_\kappa^!$
sees as negative curvature $-\kappa \cdot \omega_g$.
\end{remark}

\subsubsection{The platonic chain and the sewing hierarchy}

\begin{remark}[The platonic chain for Fredholm partition functions]%
\label{rem:thqg-X-platonic-chain}%
\index{platonic ideal form!Fredholm partition functions}%
The platonic chain of Remark~\ref{rem:platonic-ideal-form}
has the following concrete incarnation for Fredholm partition
functions:
\[
\underbrace{\algcore}_{\text{algebraic OPE data}}
\;\subset\;
\underbrace{\sewenv}_{\text{sewing envelope}}
\;\rightsquigarrow\;
\underbrace{\operatorname{Sym} A^2(D)}_{\text{Bergman Fock space}}
\;\rightsquigarrow\;
\underbrace{Z_g = \det(1 - \sewker_g)^{-\kappa}}_{\text{Fredholm
partition function}}.
\]
The first arrow (dense embedding) is the mode--Bergman
correspondence; the second arrow (Hilbert completion) uses
the sewing envelope topology; the third arrow (Fredholm
determinant) uses the one-particle reduction and the
Polyakov formula.
Each stage is necessary:
\begin{itemize}
\item $\algcore$ alone gives only formal power series in~$q$
  (the characters, which converge for $|q| < 1$ but have no
  intrinsic analytic continuation).
\item $\sewenv$ promotes the formal power series to operators
  on a Hilbert space, with explicit spectral data.
\item The Fredholm determinant provides the non-perturbative
  completion: it defines $Z_g$ as a function on $\Siegel_g$
  (the Siegel upper half-space), not merely as a formal
  $q$-series.
\end{itemize}
\end{remark}

\subsubsection{Open problems and the analytic frontier}

\begin{openproblem}[Lattice sewing envelope]%
\label{prob:thqg-X-lattice-sewing}%
\index{lattice algebra!sewing envelope}%
Construct the sewing envelope of a lattice vertex algebra
$V_\Lambda$ explicitly as a completion of the Bergman Fock
space tensored with the lattice group algebra.
Theorem~\textup{\ref{thm:lattice-sewing}} gives
\begin{equation}\label{eq:thqg-X-lattice-sewing-conj}
V_\Lambda^{\mathrm{sew}}
\;\cong\;
\operatorname{Sym} A^2(D)^{\otimes r}
\;\otimes\;
\ell^2(\Lambda^*/\Lambda),
\end{equation}
where $r = \mathrm{rank}(\Lambda)$ and
$\ell^2(\Lambda^*/\Lambda)$ is the Hilbert space of the
charge lattice.
\end{openproblem}

\begin{openproblem}[Boundary bar duality]%
\label{prob:thqg-X-boundary-bar}%
\index{bar coalgebra!boundary duality}%
Establish a Fredholm-level analytic Koszul duality:
for every analytic Koszul pair $(\cA, \cA^!)$ satisfying
HS-sewing, construct a Fredholm equivalence
\begin{equation}\label{eq:thqg-X-boundary-bar-conj}
\codercat(\anbar(\cA)\text{-comod})
\;\simeq\;
\contrcat(\cA^!\text{-mod})
\end{equation}
between the coderived category of analytic bar comodules and
the contraderived category of $\cA^!$-modules.
This would be the analytic avatar of bar-cobar inversion
(Theorem~B) at genus~$g \geq 1$.
\end{openproblem}

\begin{openproblem}[Non-Fredholm corrections: closed-form expressions]%
\label{prob:thqg-X-non-fredholm-closed}%
\index{non-Fredholm correction!closed-form}%
For class-L algebras (e.g.\ $V_k(\fg)$), express the
genus-$g$ non-Fredholm correction $\delta Z_g^{\mathrm{L}}$
in terms of Siegel modular forms and conformal blocks.
At integrable level, the Verlinde formula
\eqref{eq:verlinde-general} provides an independent
computation; the problem is to derive the Verlinde formula
directly from the Feynman integral expansion
\eqref{eq:thqg-X-class-L-correction}.
\end{openproblem}

\begin{openproblem}[Analytic completion for W-algebras]%
\label{prob:thqg-X-W-algebra-analytic}%
\index{W-algebra!analytic completion}%
Construct the sewing envelope and analytic bar coalgebra for
class-M algebras ($\mathrm{Vir}_c$, $\cW_N$) in the
resonance-filtered topology of
Theorem~\ref{thm:resonance-filtered-bar-cobar}.
The resonance rank $\rho(\cA)$
(Definition~\ref{def:resonance-rank}) controls the
difficulty: for $\rho < \infty$ (the resonance completion
theorem, Theorem~\ref{thm:platonic-completion}), the
analytic completion reduces to a finite resonance problem;
for $\rho = \infty$ (if such algebras exist among modular
Koszul chiral algebras), the completion requires a genuinely
infinite-dimensional analytic framework beyond the Fredholm
theory of this section.
\end{openproblem}

\begin{remark}[Summary: the Fredholm partition function as
a projection of $\Theta_\cA$]%
\label{rem:thqg-X-fredholm-summary}%
\index{Fredholm determinant!as MC projection}%
The genus-$g$ Fredholm partition function
$Z_g(\cA;\,\Omega) = \det(1 - \sewker_g)^{-\kappa(\cA)}$
is the analytic realization of the scalar genus tower
$F_g(\cA) = \kappa(\cA) \cdot \lambda_g^{\mathrm{FP}}$
of Theorem~D.  It is the bottom projection of the single
universal MC element
$\Theta_\cA \in \MC(\gAmod)$
through the analytic completion:
\[
\Theta_\cA
\;\xrightarrow{\;\mathrm{arity}\;2\;}
\kappa(\cA)
\;\xrightarrow{\;\mathrm{sewing}\;}
\sewker_g
\;\xrightarrow{\;\det\;}
Z_g.
\]
For class-G algebras, this projection is lossless: $Z_g$
captures all the information in~$\Theta_\cA$.  For classes
L, C, M, the non-Fredholm corrections $\delta Z_g$ carry the
information from the higher-arity projections of~$\Theta_\cA$
(the cubic shadow $\fC$, the quartic resonance class
$\mathfrak{Q}$, and all higher obstruction classes $o_r$).
The full partition function
$Z_g^{\mathrm{full}} = Z_g^{\mathrm{G}} + \delta Z_g$
is the analytic realization of the complete MC element, and
the Feynman expansion \eqref{eq:thqg-X-feynman-expansion}
is the graph-sum formula of the shadow obstruction tower
evaluated on the analytic sewing data.

This completes the development of Result~\textbf{(G10)}: the
genus-$g$ partition function of a modular Koszul chiral algebra,
computed as a nonperturbative Fredholm determinant on the
one-particle Bergman space, with explicit Feynman corrections
for non-Gaussian shadow obstruction towers.
\end{remark}

% =====================================================================
%
%  SECTION 7: GRAVITATIONAL COMPLEMENTARITY AND QUANTUM ENTROPY
%
% =====================================================================

\subsection{Gravitational complementarity and quantum
black hole entropy}
\label{subsec:thqg-X-gravitational-complementarity}
\index{gravitational entropy|textbf}%
\index{complementarity!gravitational|textbf}%
\index{BTZ black hole!quantum corrections}%
\index{CohFT!shadow|textbf}%

The preceding sections realize the scalar projection of
$\Theta_\cA$ as a Fredholm determinant.  We now develop
the \emph{full CohFT-level} consequences for gravitational
partition functions.  The derivation pipeline is:
\[
\Theta_\cA \in \MC(\gAmod)
\;\xrightarrow{\;\text{project}\;}
\Omega_{g,n}(\cA) \in H^*(\bar{\cM}_{g,n})
\;\xrightarrow{\;\text{graph sum}\;}
F_g(\cA) = \int_{\bar{\cM}_g}\!\Omega_{g,0}
\;\xrightarrow{\;\text{Legendre}\;}
S_{\mathrm{grav}}(E;\,c)
\]
where the shadow CohFT class $\Omega_{g,n}(\cA)$
is determined by the Givental $R$-matrix
(Theorem~\ref{thm:cohft-reconstruction}) and
the genus-$0$ shadow data, and the gravitational
entropy is obtained by Legendre transform of the
genus expansion.  Complementarity (Theorem~C) acts
at \emph{every stage}: it constrains the CohFT classes
at dual central charges, forces the free energies to
satisfy the Lagrangian splitting identity, and couples
the non-perturbative instanton spectra via the spectral
curve involution.

% ------------------------------------------------------------------
\subsubsection{The derivation pipeline:
from MC element to CohFT vertex factors}
\label{subsubsec:derivation-pipeline}

The shadow CohFT (Theorem~\ref{thm:shadow-cohft}) assigns
to each chirally Koszul algebra~$\cA$ a collection of
tautological classes
$\Omega_{g,n}(\cA) \in H^*(\bar{\cM}_{g,n},\Q)$
satisfying the CohFT axioms (equivariance under $S_n$,
gluing at separating and non-separating nodes).

\begin{definition}[CohFT vertex factor]%
\label{def:cohft-vertex-factor}%
\index{CohFT!vertex factor}%
For integers $g \geq 0$, $n \geq 0$ with $2g-2+n > 0$,
the \emph{CohFT vertex factor} of~$\cA$ at $(g,n)$ is
\begin{equation}\label{eq:cohft-vertex-def}
V(g,n;\,\cA)
\;:=\;
\int_{\bar{\cM}_{g,n}}
\Omega_{g,n}(\cA).
\end{equation}
At genus~$0$: $V(0,n;\,\cA) = S_n(\cA)$, the genus-$0$
shadow coefficient at arity~$n$
($S_2 = \kappa$, $S_3 = \alpha$, $S_4 = Q^{\mathrm{contact}}$, etc.).
These are \emph{family-specific} and encode the OPE structure.
\end{definition}

\begin{proposition}[Givental reconstruction of
higher-genus vertex factors]%
\label{prop:givental-vertex-reconstruction}%
\ClaimStatusProvedHere%
\index{Givental reconstruction!vertex factors}%
For $g \geq 1$ and $n \geq 0$ with $2g-2+n > 0$,
the vertex factor~\eqref{eq:cohft-vertex-def} is
determined by the Givental $R$-matrix
$R_\cA(z) = \sum_{d \geq 0} R_d\, z^d$
$($Theorem~\ref{thm:cohft-reconstruction}$)$
and the Witten--Kontsevich intersection numbers:
\begin{equation}\label{eq:givental-vertex-formula}
V(g,n;\,\cA)
\;=\;
\sum_{\substack{d_1,\ldots,d_n \geq 0 \\[2pt]
d_1+\cdots+d_n = 3g-3+n}}
\Bigl(\prod_{i=1}^n R_{d_i}\Bigr)
\cdot
\langle\tau_{d_1}\cdots\tau_{d_n}\rangle_g
\end{equation}
where
$\langle\tau_{d_1}\cdots\tau_{d_n}\rangle_g
= \int_{\bar{\cM}_{g,n}}
\psi_1^{d_1}\cdots\psi_n^{d_n}$
are computed by the DVV recursion.

The $R$-matrix is the complementarity propagator
$($Theorem~\ref{thm:cohft-reconstruction}$)$:
for the Hodge sector,
$R(z) = \exp\!\bigl(\sum_{k \geq 1}
\tfrac{B_{2k}}{2k(2k-1)}\,z^{2k-1}\bigr)$,
giving
$R_0 = 1$,
$R_1 = \tfrac{1}{12}$,
$R_2 = \tfrac{1}{288}$.
The $R$-matrix is \textbf{universal}:
it depends only on the Hodge structure of
$\bar{\cM}_{g,n}$, not on the algebra~$\cA$.
Family-specific information enters through the
genus-$0$ vertex factors $V(0,n;\,\cA) = S_n$.
\end{proposition}

\begin{proof}
This is the rank-$1$ specialization of the
Givental--Teleman reconstruction theorem
(Theorem~\ref{thm:cohft-reconstruction}):
the shadow CohFT is obtained from the trivial CohFT
by the quantized $R$-matrix action
$\hat{R} \cdot \Omega^{\mathrm{triv}}$,
where $R = P_\cA$ is the complementarity propagator.
The $R(z)$-insertion at each leg of a vertex
produces the $\psi$-class weights $R_{d_i}$, and
integration over $\bar{\cM}_{g,n}$ yields the
Witten--Kontsevich numbers by definition.
\end{proof}

\begin{remark}[Hodge-lambda integral vs.\ CohFT
partition function]%
\label{rem:hodge-lambda-distinction}%
\index{Hodge-lambda integral}%
The genus-$g$ free energy from Theorem~D is:
\begin{equation}\label{eq:fg-hodge-lambda}
F_g^{\mathrm{sc}}(\cA)
\;=\;
\kappa(\cA)\,\lambda_g^{\mathrm{FP}}
\;=\;
\kappa(\cA)
\int_{\bar{\cM}_{g,1}}
\lambda_g\,\psi^{2g-2},
\end{equation}
a Hodge-lambda integral on $\bar{\cM}_{g,1}$
$($one marked point, with $\psi$-insertion$)$.
The formula~\eqref{eq:givental-vertex-formula}
computes $\psi$-class integrals
$\int \prod R(\psi_i)$ \emph{without} the
Hodge class~$\lambda_g$.  These are structurally
different: the shadow obstruction tower integrand involves
$\lambda_g$ $($from the Hodge bundle in the modular
operad's Feynman transform$)$, while the CohFT
vertex factor integrates only $\psi$-classes
weighted by~$R$.  Shadow corrections
modify the class $\lambda_g$ by shadow obstruction tower data,
\emph{not} the propagator.
\end{remark}

\begin{remark}[Modular operad vertex factors from the MC equation]%
\label{rem:operad-vertex-factors}%
\index{modular operad!vertex factors|textbf}%
The MC equation $D_\cA^2 = 0$, decomposed at each
$(g,n)$, determines the modular operad vertex factors
$V^{\mathrm{op}}(g,n;\,\cA)$ from lower-level shadow data.
At $(1,1)$, the genus-loop of the cubic shadow gives:
\begin{equation}\label{eq:V11-mc}
V^{\mathrm{op}}(1,1;\,\cA)
\;=\;
-\frac{3\,S_3}{\kappa}
\;=\;
-\frac{3\,\alpha}{\kappa}
\end{equation}
$($vanishes for Gaussian; equals $-12/c$ for Virasoro$)$.
At $(1,2)$, combining the cubic self-sewing and the
genus-loop of the quartic:
\begin{equation}\label{eq:V12-mc}
V^{\mathrm{op}}(1,2;\,\cA)
\;=\;
\frac{3\,S_3^2}{\kappa^2}
\;-\;
\frac{6\,S_4}{\kappa}
\end{equation}
$($vanishes for Gaussian; equals
$(240c{+}936)/(c^2(5c{+}22))$ for Virasoro$)$.
These are the \emph{modular operad} vertex factors,
distinct from the $\psi$-class vertex factors $T(g,n)$
of~\eqref{eq:givental-vertex-formula}
$($Remark~\ref{rem:hodge-lambda-distinction}$)$.
\end{remark}

% ------------------------------------------------------------------
\subsubsection{The gravitational genus expansion}
\label{subsubsec:gravitational-genus-expansion}

\begin{theorem}[Gravitational complementarity genus expansion]%
\label{thm:gravitational-complementarity-genus-expansion}%
\ClaimStatusProvedHere%
\index{complementarity!genus expansion}%
Let $c > 0$ with $c \neq 26$.
Write $\kappa = \kappa(\mathrm{Vir}_c) = c/2$ and define:
\begin{itemize}
\item the \textbf{scalar free energy}
$F_g^{\mathrm{sc}}(c) := \kappa\,\lambda_g^{\mathrm{FP}}$
where
$\lambda_g^{\mathrm{FP}}
= (2^{2g-1}-1)|B_{2g}|/(2^{2g-1}(2g)!)$;
\item the \textbf{full free energy}
$F_g(c) := \sum_{\Gamma}
|\operatorname{Aut}(\Gamma)|^{-1}\,w_\Gamma$,
the graph sum over stable graphs of genus~$g$
with $n = 0$ marked points, where $w_\Gamma$
uses the CohFT vertex factors from
Proposition~\ref{prop:givental-vertex-reconstruction}
at higher-genus vertices and $V(0,n) = S_n$
at genus-$0$ vertices;
\item the \textbf{shadow correction}
$\delta F_g(c) := F_g(c) - F_g^{\mathrm{sc}}(c)$,
which receives contributions only from graphs containing
at least one genus-$0$ vertex of valence $\geq 3$
$($i.e., using the shadow data $S_3, S_4, \ldots)$.
\end{itemize}

\begin{enumerate}[label=\textup{(\roman*)}]
\item \textbf{Scalar complementarity.}
For every $g \geq 1$:
\begin{equation}\label{eq:scalar-complementarity-genus}
F_g^{\mathrm{sc}}(c) + F_g^{\mathrm{sc}}(26-c)
\;=\;
13\,\lambda_g^{\mathrm{FP}}.
\end{equation}

\item \textbf{$\hat{A}$-genus generating function.}
The scalar generating function has radius of
convergence~$2\pi$ in~$\hbar$:
\begin{equation}\label{eq:ahat-genus-gf}
\sum_{g=1}^{\infty}
F_g^{\mathrm{sc}}(c)\,\hbar^{2g-2}
\;=\;
\frac{c}{2\hbar^2}
\Bigl[\hat{A}(i\hbar) - 1\Bigr],
\end{equation}
where
$\hat{A}(x) = (x/2)/\sinh(x/2)$
is the $\hat{A}$-genus and
$\hat{A}(ix) = (x/2)/\sin(x/2)$
has all positive Taylor coefficients, matching
$F_g^{\mathrm{sc}} > 0$.  The first pole of
$(x/2)/\sin(x/2)$ at $x = 2\pi$ gives the stated
radius of convergence.

\item \textbf{Shadow correction Borel structure.}
The shadow corrections satisfy
$\delta F_g \sim C(c)\,\rho(c)^g\,g^{-5/2}$
as $g \to \infty$,
where $\rho(c) = \sqrt{(180c+872)/(c^2(5c+22))}$ is the
shadow growth rate
$($Theorem~\ref{thm:shadow-radius}$)$.
The correction series has:
\begin{itemize}
\item convergence radius
$R = 2\pi/\rho(c) > 2\pi$ for
$c > c^* \approx 6.125$;
\item Gevrey-$1$ divergence, Borel summable along
$\arg(\hbar^2) \neq 0$, for $c < c^*$.
\end{itemize}
The scalar series converges for $|\hbar| < 2\pi$
$($the first pole of $\hat{A}(i\hbar) = (\hbar/2)/\sin(\hbar/2)$ at $\hbar = 2\pi$$)$;
only the shadow corrections can diverge \emph{faster}.

\item \textbf{Self-dual point.}
At $c = 13$:
$\mathrm{Vir}_{13}^! = \mathrm{Vir}_{13}$,
$\Delta(13) = \Delta(13)^! = 40/87$,
$\rho(13) \approx 0.467 < 1$.
Complementarity forces
$2F_g(13) = \Phi_g(13)$,
pinning down the self-dual free energies from
$\bar{\cM}_g$ topology alone.
\end{enumerate}
\end{theorem}

\begin{proof}
\textup{(i)} From $\kappa + \kappa^! = c/2 + (26-c)/2 = 13$.
\textup{(ii)} From the Taylor expansion of
$\hat{A}(ix) = (x/2)/\sin(x/2)
= 1 + \sum_{g \geq 1}\lambda_g^{\mathrm{FP}}\,x^{2g}$
(all coefficients positive) and
$F_g^{\mathrm{sc}} = (c/2)\lambda_g^{\mathrm{FP}}$.
\textup{(iii)} By
Theorem~\ref{thm:shadow-radius}, the $S_r$ grow as
$S_r \sim C\rho^r r^{-5/2}$ for class-M algebras;
this controls the graph-sum asymptotics at large~$g$
via the dominant contribution from maximally trivalent
stable graphs.  The scalar series
$(c/2\hbar^2)(\hat{A}(i\hbar) - 1)$ has
radius of convergence~$2\pi$, from the first pole of
$(\hbar/2)/\sin(\hbar/2)$ at $\hbar = 2\pi$.
\textup{(iv)}
$\kappa(13) = 13/2 = \kappa(13)^!$;
$\Delta(13) = 8\kappa S_4|_{c=13}
= 40/(5\cdot 13+22) = 40/87$;
$\rho(13)^2 = (180\cdot 13+872)/(169\cdot 87)
= 3212/14703 \approx 0.218$.
\end{proof}

% ------------------------------------------------------------------
\subsubsection{Genus-$2$ free energy with shadow corrections}
\label{subsubsec:genus2-shadow-corrections}

\begin{proposition}[Full genus-$2$ free energy for Virasoro]%
\label{prop:virasoro-F2-full}%
\ClaimStatusProvedHere%
\index{free energy!genus 2!Virasoro|textbf}%
Using the MC-derived vertex
factors~\eqref{eq:V11-mc}--\eqref{eq:V12-mc} and the
$6$ stable graphs at $(g{=}2, n{=}0)$, the genus-$2$ free
energy of $\mathrm{Vir}_c$ is
\begin{equation}\label{eq:virasoro-F2-full}
F_2(\mathrm{Vir}_c)
\;=\;
\underbrace{\frac{7c}{11520}}_{\text{scalar}}
\;+\;
\underbrace{
\frac{240c+941}{c^3(5c+22)}
\;+\;
\frac{296}{3c^3}
}_{\text{shadow corrections}}.
\end{equation}
The shadow corrections come from four non-corolla graphs:
the theta graph $(\alpha^2 P^3/12)$, the banana
$(S_4\,P^2/8)$, the separating node
$((V^{\mathrm{op}}_{1,1})^2 P/2)$,
and the irreducible node
$(V^{\mathrm{op}}_{1,2}\,P/2)$,
minus the mixed graph
$(\alpha\,V^{\mathrm{op}}_{1,1}\,P^2/2)$.
At $c = 13$ $($self-dual$)$:
$F_2(13) \approx 0.074$.
\end{proposition}

\begin{proof}
Direct computation from the graph sum with
$V(0,2) = c/2$,
$V(0,3) = 2$,
$V(0,4) = 10/(c(5c+22))$,
$V^{\mathrm{op}}(1,1) = -12/c$
(equation~\eqref{eq:V11-mc}),
$V^{\mathrm{op}}(1,2) = (240c+936)/(c^2(5c+22))$
(equation~\eqref{eq:V12-mc}),
$V(2,0) = 7c/11520$,
and propagator $P = 2/c$.
Verified computationally: $F_2(13) = F_2(13)^!$
(self-duality at $c = 13$), confirming
complementarity at the non-scalar level.
\end{proof}

% ------------------------------------------------------------------
\subsubsection{CohFT-level complementarity and Lagrangian splitting}
\label{subsubsec:cohft-lagrangian}

\begin{proposition}[CohFT complementarity]%
\label{prop:cohft-complementarity}%
\ClaimStatusProvedHere%
\index{complementarity!CohFT-level}%
The shadow CohFT satisfies the Lagrangian complementarity
identity at the \emph{class level}: for every
$g \geq 1$ and $n \geq 0$ with $2g-2+n > 0$,
\begin{equation}\label{eq:cohft-complementarity}
\Omega_{g,n}(\cA) \oplus \Omega_{g,n}(\cA^!)
\;\hookrightarrow\;
H^*(\bar{\cM}_{g,n},\,\cZ(\cA))
\end{equation}
where $\cZ(\cA)$ is the center sheaf,
and the inclusion is a Lagrangian embedding
with respect to the $({-}(3g{-}3{+}n))$-shifted
symplectic structure from
Theorem~\ref{thm:ambient-complementarity-fmp}.
Integrating~\eqref{eq:cohft-complementarity}
over $\bar{\cM}_{g,n}$ yields the
\emph{vertex-factor complementarity}:
\begin{equation}\label{eq:vertex-complementarity}
V(g,n;\,\cA) + V(g,n;\,\cA^!)
\;=\;
\Phi_{g,n}(\cA)
\end{equation}
where $\Phi_{g,n}$ is determined by the center sheaf.
At the scalar level:
$\Phi_{g,0} = \kappa \lambda_g^{\mathrm{FP}}
+ \kappa^!\lambda_g^{\mathrm{FP}}
= 13\,\lambda_g^{\mathrm{FP}}$
for Virasoro, recovering~\textup{(i)}.
\end{proposition}

\begin{proof}
Apply Theorem~C (Complementarity)
at the CohFT level: the decomposition
$\mathbf{Q}_g(\cA) \oplus \mathbf{Q}_g(\cA^!)
\simeq H^*(\bar{\cM}_g, \cZ(\cA))$
is a Lagrangian splitting of the shifted symplectic
structure on $H^*(\bar{\cM}_g, \cZ(\cA))$, where
$\cZ(\cA)$ carries the perfect pairing from
Verdier duality.  The CohFT class
$\Omega_{g,n}(\cA) \in \mathbf{Q}_g(\cA)$ is
the $\cA$-half of this splitting.
Integration over $\bar{\cM}_{g,n}$ projects
the Lagrangian identity to the vertex-factor
identity~\eqref{eq:vertex-complementarity}.
\end{proof}

% ------------------------------------------------------------------
\subsubsection{Quantum black hole entropy}
\label{subsubsec:quantum-bh-entropy}

\begin{remark}[BTZ entropy from the genus expansion]%
\label{rem:quantum-bh-entropy}%
\index{Bekenstein--Hawking entropy!quantum corrections}%
For 3D gravity with Brown--Henneaux central charge
$c = 3\ell/(2G_N)$, the genus expansion of the
gravitational partition function yields the BTZ entropy
with quantum corrections:
\begin{equation}\label{eq:btz-entropy-expansion}
S_{\mathrm{BTZ}}(E;\,c)
\;=\;
2\pi\sqrt{\frac{cE}{6}}
\;-\;
\tfrac{3}{2}\,\log\frac{cE}{6}
\;+\;
\sum_{g=2}^{\infty}
\alpha_g(c)\,
\Bigl(\frac{6}{cE}\Bigr)^{\!g-1}
\end{equation}
where the leading term is the Cardy/Bekenstein--Hawking
entropy $($genus~$0$, saddle-point of the modular
$S$-transform$)$, the logarithmic correction
$-\tfrac{3}{2}\log(cE/6)$ is the one-loop
determinant $($genus~$1)$, and the $\alpha_g(c)$ are
computable from the genus-$g$ free energies~$F_g(c)$
via saddle-point expansion.

The shadow corrections $\delta F_g$ give the
first \emph{systematically computable} quantum gravity
corrections beyond the logarithmic term.
At the scalar level,
$\alpha_g^{\mathrm{sc}} = F_g^{\mathrm{sc}}(c)
= (c/2)\lambda_g^{\mathrm{FP}}$.

\textbf{Complementarity constrains the corrections
at dual central charges.}
The Lagrangian identity
$F_g(c) + F_g(26-c) = \Phi_g(c)$ means the
black hole entropy at central charge~$c$ and at~$26-c$
are \emph{not independent}: knowing one determines
the other through the center sheaf.
At $c = 13$, complementarity pins down the quantum
corrections from $\bar{\cM}_g$ topology alone.
\end{remark}

% ------------------------------------------------------------------
\subsubsection{Instanton structure and spectral curve}
\label{subsubsec:instanton-spectral}

\begin{remark}[Discriminant complementarity and
instanton duality]%
\label{rem:discriminant-instanton-duality}%
\index{discriminant!complementarity}%
\index{instanton!complementarity}%
The critical discriminants satisfy
$($Theorem~\ref{thm:shadow-connection}$)$:
\begin{equation}\label{eq:discriminant-complementarity-qg}
\Delta(\mathrm{Vir}_c) + \Delta(\mathrm{Vir}_{26-c})
\;=\;
\frac{6960}{(5c+22)(152-5c)}.
\end{equation}
The spectral curve $y^2 = Q_L(t)$ of the shadow
tower~$($Corollary~\ref{cor:spectral-curve}$)$
determines the Borel structure of the genus expansion:
\begin{itemize}
\item its \emph{branch points} are the instanton
locations in the Borel plane, at distance
$1/\rho(c)$ from the origin;
\item its \emph{periods} give the instanton actions
$A(c) = 1/\rho(c)$;
\item its \emph{monodromy} $= -1$
is the Koszul involution
(Theorem~\ref{thm:shadow-connection}).
\end{itemize}
The discriminant
sum~\eqref{eq:discriminant-complementarity-qg}
constrains the pair $(\rho(c), \rho(26-c))$ through
the shadow metric.  At $c = 13$, the spectral curve
has enhanced $\Z_2$ symmetry
$(\Delta = \Delta^! = 40/87)$:
the instanton spectrum is self-dual.

The \emph{Eynard--Orantin topological recursion}
on the spectral curve $y^2 = Q_L(t)$
$($Corollary~\ref{cor:topological-recursion-mc-shadow}$)$
computes the full genus-$g$ free energy $F_g(\cA)$
including \emph{all} shadow corrections, via the
residue formula at the branch points.
For the Virasoro spectral curve
$y^2 = c^2 + 12ct + \frac{180c+872}{5c+22}\,t^2$,
this gives $F_g(\mathrm{Vir}_c)$ as an explicit rational
function of~$c$ at each genus.
\end{remark}
